\chapter{Formation of Singularities}
\label{chap:cz-singularity-formation}

\bigskip
Let $g_{ij}(x,t)$ be a solution to the Ricci flow on $M\times
[0,T)$ and suppose $[0,T)$, $T\le \infty$, is the maximal time
interval. If $T<+\infty$, then the short time existence theorems
in Section 1.2 tells us the curvature of the solution becomes
unbounded as $t\rightarrow T$ (cf. Theorem 8.1 of \cite{Ha95F}).
We then say the solution \textbf{develops a
singularity}\index{solution develops a singularity} as
$t\rightarrow T$. As in the minimal surface theory and harmonic
map theory, one usually tries to understand the structure of a
singularity of the Ricci flow by rescaling the solution (or blow
up) to obtain a sequence of solutions to the Ricci flow with
uniformly bounded curvature on compact subsets and looking at its
limit.

The main purpose of this chapter is to present the convergence
theorem of Hamilton \cite{Ha95} for a sequence of solutions to the
Ricci flow with uniform bounded curvature on compact subsets, and
to use the convergence theorem to give a rough classification in
\cite{Ha95F} for singularities of solutions to the Ricci flow.
Further studies on the structures of singularities of the Ricci
flow will be given in Chapter 6 and 7.


\section{Cheeger Type Compactness}

In this section, we establish Hamilton's compactness theorem
(Theorem 4.1.5) for solutions to the Ricci flow. The presentation
is based on Hamilton \cite{Ha95}.

We begin with the concept of $C_{loc}^{\infty}$ convergence of
tensors on a given manifold $M$. Let ${T_i}$ be a sequence of
tensors on $M$. We say that \mbox{\boldmath{$T_i$}} {\bf converges
to a tensor} \mbox{\boldmath{$T$}}\index{$T_i$ converges to a
tensor $T$} in the $C_{loc}^{\infty}$ topology if we can find a
covering $\{(U_s,\varphi_s)\}$, $\varphi_s: U_s\rightarrow
\mathbb{R}^n$, of $C^\infty$ coordinate charts so that for every
compact set $K\subset M$, the components of $T_i$ converge in the
$C^\infty$ topology to the components of $T$ in the intersections
of $K$ with these coordinate charts, considered as functions on
$\varphi_s(U_s)\subset \mathbb{R}^n$. Consider a Riemannian
manifold $(M,g)$. A {\bf marking}\index{marking} on $M$ is a
choice of a point $p\in M$ which we call the {\bf
origin}\index{origin}. We will refer to such a triple $(M,g,p)$ as
a {\bf marked Riemannian manifold}.\index{marked Riemannian
manifold}

%\vskip0.1cm \noindent 4.1.1
% CZ-BASELINE source-line:6894 source-kind:definition source-id:4.1.1
\begin{definition}[pointed smooth convergence of Riemannian manifolds]
  \label{def:cz-pointed-smooth-convergence-of-riemannian-manifolds}
  \dcref{Definition 4.1.1}
  \group{Cao--Zhu Cheeger Type Compactness}
  \source{cao-zhu}
Let $(M_k,g_k,p_k)$ be a sequence of marked complete Riemannian
manifolds, with metrics $g_k$ and marked points $p_k\in M_k$. Let
$B(p_k, s_k)\subset M_k$ denote the geodesic ball centered at
$p_k\in M_k$ and of radius $s_k$ ($0 < s_k \leq +\infty$). We say
a sequence of marked geodesic balls $(B(p_k,s_k),g_k, p_k)$ with
$s_k \rightarrow s_{\infty}(\leq +\infty)$ {\bf converges} in the
$C_{loc}^{\infty}$ topology {\bf to a marked} (maybe noncomplete)
{\bf manifold}\index{converges to a marked manifold}
$(B_{\infty},g_{\infty},p_{\infty})$, which is an open geodesic
ball centered at $p_{\infty}\in B_{\infty}$ and of radius
$s_{\infty}$ with respect to the metric $g_{\infty}$, if we can
find a sequence of exhausting open sets $U_k$ in $B_{\infty}$
containing $p_{\infty}$ and a sequence of diffeomorphisms $f_k$ of
the sets $U_k$ in $B_{\infty}$ to open sets $V_k$ in $B(p_k,s_k)
\subset M_k$ mapping $p_{\infty}$ to $p_k$ such that the pull-back
metrics $\tilde{g}_k=(f_k)^*g_k$ converge in $C^\infty$ topology
to $g_{\infty}$ on every compact subset of $B_{\infty}$.
\end{definition}

We remark that this concept of $C_{\rm loc}^\infty$-convergence of
a sequence of marked manifolds $(M_k,g_k,p_k)$ is not the same as
that of $C_{\rm loc}^\infty$-convergence of metric tensors on a
given manifold, even when we are considering the sequence of
Riemannian metric $g_k$ on the same space $M$. This is because one
can have a sequence of diffeomorphisms $f_k: M\rightarrow M$ such
that $(f_k)^*g_k$ converges in $C^\infty_{\rm loc}$ topology while
$g_k$ itself does not converge.

There have been a lot of work in Riemannian geometry on the
convergence of a sequence of compact manifolds with bounded
curvature, diameter and injectivity radius (see for example Gromov
\cite{Gro}, Peters \cite{Pe}, and Greene and Wu \cite{Gre}). The
following theorem, which is a slight generalization of Hamilton's
convergence theorem \cite{Ha95}, modifies these results in three
aspects: the first one is to allow noncompact limits and then to
avoid any diameter bound; the second one is to avoid having to
assume a uniform lower bound for the injectivity radius over the
whole manifold, a hypothesis which is much harder to satisfy in
applications; the last one is to avoid a uniform curvature bound
over the whole manifold so that we can take a local limit.

%\vskip 0.2cm\noindent 4.1.2
% CZ-BASELINE source-line:6937 source-kind:theorem source-id:4.1.2
\begin{theorem}[Cheeger-type pointed compactness]
  \label{thm:cz-cheeger-type-pointed-compactness}
  \dcref{Theorem 4.1.2}
  \group{Cao--Zhu Cheeger Type Compactness}
  \source{cao-zhu}
Let\; $(M_k,\,g_k,\,p_k)$\; be\; a\; sequence\; of marked complete
Riemannian manifolds of dimension $n$. Consider a sequence of
geodesic balls $B(p_k,s_k) \subset M_k$ of radius $s_k$ $(0<s_k\le
\infty)$, with $s_k \rightarrow s_{\infty}(\le \infty)$, around
the base point $p_k$ of $M_k$ in the metric $g_k$. Suppose
\begin{itemize}
\item[(a)] for every radius $r<s_{\infty}$ and every integer $l\ge
0$ there exists a constant $B_{l,r}$, independent of $k$, and
positive integer $k(r,l) < +\infty$ such that as $k \geq k(r,l)$,
the curvature tensors $Rm(g_k)$ of the metrics $g_k$ and their
$l^{\mbox{th}}$-covariant derivatives satisfy the bounds
$$
|\nabla^lRm(g_k)|\leq B_{l,r}
$$
on the balls $B(p_k, r)$ of radius $r$ around $p_k$ in the metrics
$g_k$; and \item[(b)] there exists a constant $\delta>0$
independent of $k$ such that the injectivity radii
inj$\,(M_k,p_k)$ of $M_k$ at $p_k$ in the metric $g_k$ satisfy the
bound
$$
{\rm inj}\,(M_k,p_k)\geq\delta.$$
\end{itemize}
Then there exists a subsequence of the marked geodesic balls
$(B(p_k,\,s_k),\,g_k,$ $p_k)$ which converges to a marked geodesic
ball $(B(p_{\infty},s_{\infty}),g_{\infty},p_{\infty})$ in
$C^\infty_{\rm loc}$ topology. Moreover the limit is complete if
$s_{\infty} = +\infty$.

\end{theorem}

%\vskip 0.3cm \noindent {\bf Proo.}
\begin{proof}
  \uses{cor:cz-injectivity-radius-from-local-curvature-scale,prop:cz-local-metric-subsequence-convergence,prop:cz-transition-isometry-subsequence-convergence}
In \cite{Ha95} (see also Theorem 16.1 of \cite{Ha95F}), Hamilton
proved the above convergence theorem for the case $s_{\infty} =
+\infty$. Thus it remains to prove the case of $s_{\infty} <
+\infty$. Our proof here follows, with some slight modifications,
the argument of Hamilton \cite{Ha95}. In fact, except Step 1, the
proof is essentially taken from Hamilton \cite{Ha95}. Suppose we
are given a sequence of geodesic balls $(B(p_k,s_k),g_k, p_k)
\subset (M_k,g_k,p_k)$, with $s_k \rightarrow s_{\infty} (<
+\infty)$, satisfying the assumptions of Theorem 4.1.2. We will
split the proof into three steps.

\medskip
{\it Step} 1: \ Picking the subsequence.

\smallskip
By the local injectivity radius estimate (4.2.2) in Corollary
4.2.3 of the next section, we can find a positive decreasing $C^1$
function $\rho (r)$, $0 \leq r < s_{\infty}$, independent of $k$
such that
\begin{align}
\rho(r) &< \frac{1}{100}(s_{\infty} - r), \\  %\eqno (4.1.1)
0 &\geq \rho'(r) \geq -\frac{1}{1000}, %%\eqno (4.1.2)
\end{align}
and a sequence of positive constants $\varepsilon_k \rightarrow 0$
so that the injectivity radius at any point $x\in B(p_k,s_k)$ with
$r_k = d_k(x,p_k) < s_{\infty} - \varepsilon_k$ is bounded from
below by \be
{\rm inj}\, (M_k, x)\ge 500\rho (r_k(x)),  %\eqno (4.1.3)
\ee where $r_k(x)=d_k(x, p_k)$ is the distance from $x$ to $p_k$
in the metric $g_k$ of $M_k$. We define
$$
\tilde\rho(r)=\rho(r+20\rho(r)),\quad
\tilde{\tilde\rho}(r)=\tilde\rho(r+20\tilde\rho(r)).
$$
By (4.1.2) we know that both $\tilde{\rho}(r)$ and
$\tilde{\tilde{\rho}}(r)$ are nonincreasing positive functions on
$[0,s_{\infty})$.

In each $B(p_k,s_{\infty})$ we choose inductively a sequence of
points $x_k^{\alpha}$ for $\alpha=0,1,2,\ldots$ in the following
way. First we let $x_k^0=p_k$. Once $x_k^{\alpha}$ are chosen for
$\alpha=0,1,2,\ldots, \sigma$, we pick $x_k^{\sigma+1}$ closest to
$p_k$ so that $r_k^{\sigma+1}=r_k(x_k^{\sigma+1})$ is as small as
possible, subject to the requirement that the open ball
$B(x_k^{\sigma+1}, \tilde{\tilde\rho}_k^{\sigma+1})$ around
$x_k^{\sigma+1}$ of radius $\tilde{\tilde\rho}_k^{\sigma+1}$ is
disjoint from the balls $B(x_k^{\alpha},
\tilde{\tilde\rho}_k^{\alpha})$ for $\alpha=0,1,2,\ldots, \sigma$,
where
$\tilde{\tilde\rho}_k^{\alpha}=\tilde{\tilde\rho}(r_k^{\alpha})$
and $r_k^{\alpha}=r_k(x_k^{\alpha})$. In particular, the open
balls $B(x_k^{\alpha}, \tilde{\tilde\rho}_k^{\alpha}),
\alpha=0,1,2,\ldots$, are all disjoint. We claim the balls
$B(x_k^{\alpha}, 2\tilde{\tilde\rho}_k^{\alpha})$ cover
$B(p_k,s_{\infty}-\varepsilon_k)$ and moreover for any $ r$, $0< r
<s_{\infty} -\varepsilon_k$, we can find $\lambda ( r)$
independent of $k$ such that for $k$ large enough, the geodesic
balls $B(x_k^{\alpha}, 2\tilde{\tilde\rho}_k^{\alpha})$ for
$\alpha\le\lambda ( r)$ cover the ball $B(p_k,r)$.

To see this, let $x\in B(p_k,s_{\infty}-\varepsilon_k)$ and let
$r(x)$ be the distance from $x$ to $p_k$ and let
$\tilde{\tilde\rho}=\tilde{\tilde\rho}(r(x))$. Consider those
$\alpha$ with $r_k^{\alpha}\le r(x) < s_{\infty} - \varepsilon_k$.
Then
$$
\tilde{\tilde\rho}\le \tilde{\tilde\rho}_k^{\alpha}.
$$
Now the given point $x$ must lie in one of the balls
$B(x_k^{\alpha}, 2\tilde{\tilde\rho}_k^{\alpha})$. If not, we
could choose the next point in the sequence of $x_k^{\beta}$ to be
$x$ instead, for since
$\tilde{\tilde\rho}_k^{\alpha}+\tilde{\tilde\rho}\le
2\tilde{\tilde\rho}_k^{\alpha}$ the ball $B(x,
\tilde{\tilde\rho})$ would miss $B(x_k^{\alpha},
\tilde{\tilde\rho}_k^{\alpha})$ with $r_k^\alpha\le r(x)$. But
this is a contradiction. Moreover for any $ r$, $0< r <s_{\infty}
-\varepsilon_k$, using the curvature bound and the injectivity
radius bound, each ball $B(x_k^{\alpha},
\tilde{\tilde\rho}_k^{\alpha})$ with $r_k^{\alpha}\le r$ has
volume at least $\epsilon(r)\tilde{\tilde\rho}^n$ where
$\epsilon(r)>0$ is some constant depending on $r$ but independent
of $k$. Now these balls are all disjoint and contained in the ball
$B(p_k, (r+s_{\infty})/2)$. On the other hand, for large enough
$k$, we can estimate the volume of this ball from above, again
using the curvature bound, by a positive function of $r$ that is
independent of $k$. Thus there is a $k'(r) > 0$ such that for each
$k \geq k'(r)$, there holds \be \# \{ \alpha \ |\ r_k^{\alpha}\le
r \} \leq \lambda(r)
%\eqno (4.1.4)
\ee for some positive constant $\lambda(r)$ depending only on $r$,
and the geodesic balls $B(x_k^{\alpha},
2\tilde{\tilde\rho}_k^{\alpha})$ for $\alpha\le\lambda ( r)$ cover
the ball $B(p_k,r)$.

By the way, since
\begin{align*}
r_k^\alpha &  \leq r_k^{\alpha -1} +
\tilde{\tilde{\rho}}_k^{\alpha -1} +\tilde{\tilde{\rho}}_k^{\alpha} \\
&  \leq r_k^{\alpha -1} + 2\tilde{\tilde{\rho}}_k^{\alpha -1},
\end{align*}
and by (4.1.1)
$$
\tilde{\tilde{\rho}}_k^{\alpha -1} \leq \frac{1}{100}(s_{\infty}
-r_k^{\alpha -1}),
$$
we get by induction
\begin{align*}
r_k^\alpha & \leq \frac{49}{50}r_k^{\alpha -1} +\frac{1}{50}s_{\infty} \\
&  \leq \(\frac{49}{50}\)^{\alpha}r_k^{0}
+\frac{1}{50}\(1+\frac{49}{50}+ \cdots
+ \(\frac{49}{50}\)^{\alpha-1}\)s_{\infty} \\
&  = \(1-\(\frac{49}{50}\)^{\alpha}\)s_{\infty}.
\end{align*}
So for each $\alpha$, with $\alpha \leq \lambda(r)$
$(r<s_{\infty})$, there holds \be r_k^\alpha \leq
\(1-\(\frac{49}{50}\)^{\lambda(r)}\)s_{\infty}
%\eqno(4.1.5)
\ee for all $k$. And by passing to a subsequence (using a
diagonalization argument) we may assume that $r_k^{\alpha}$
converges to some $r^{\alpha}$ for each ${\alpha}$. Then
$\tilde{\tilde\rho}_k^{\alpha}$ (respectively
$\tilde\rho_k^{\alpha}, \rho_k^{\alpha}$) converges to
$\tilde{\tilde\rho}^{\alpha}=\tilde{\tilde\rho}(r^{\alpha})$
(respectively $\tilde\rho^{\alpha}=\tilde\rho(r^{\alpha}),
\rho^{\alpha}=\rho(r^{\alpha})$).

Hence for all $\alpha$ we can find $k({\alpha})$ such that
$$
\frac{1}{2}\tilde{\tilde\rho}^{\alpha}
\le\tilde{\tilde\rho}_k^{\alpha} \le 2\tilde{\tilde\rho}^{\alpha}
$$
$$
\frac{1}{2}\tilde\rho^{\alpha}\le \tilde\rho_k^{\alpha}\le
2\tilde\rho^{\alpha} \quad \mbox{and} \quad
\frac{1}{2}\rho^{\alpha}\le \rho_k^{\alpha}\le 2\rho^{\alpha}
$$
whenever $k\ge k(\alpha)$. Thus for all $\alpha$,
$\tilde{\tilde\rho}_k^{\alpha}$ and $\tilde{\tilde\rho}^{\alpha}$
are comparable when $k$ is large enough so we can work with balls
of a uniform size, and the same is true for
$\tilde\rho_k^{\alpha}$ and $\tilde\rho^{\alpha}$, and
$\rho_k^{\alpha}$ and $\rho^{\alpha}$. Let $\hat
B_k^{\alpha}=B(x_k^{\alpha}, 4\tilde{\tilde\rho}^{\alpha})$, then
$\tilde{\tilde\rho}_k^{\alpha}\le 2\tilde{\tilde\rho}^{\alpha}$
and $B(x_k^{\alpha}, 2\tilde{\tilde\rho}_k^{\alpha})\subset
B(x_k^{\alpha}, 4\tilde{\tilde\rho}^{\alpha})=\hat B_k^{\alpha}$.
So for every $r$ if we let $k(r)=\max\{k(\alpha) \  |\ \alpha\le
\lambda(r)\}$ then when $k\ge k(r)$, the balls $\hat B_k^{\alpha}$
for $\alpha\le \lambda(r)$ cover the ball $B(p_k, r)$ as well.
Suppose that $\hat B_k^{\alpha}$ and $\hat B_k^{\beta}$ meet for
$k\ge k({\alpha})$ and $k\ge k({\beta})$, and suppose
$r_k^{\beta}\le r_k^{\alpha}$. Then, by the triangle inequality,
we must have
$$
r_k^{\alpha}\le r_k^{\beta}
+4\tilde{\tilde\rho}^{\alpha}+4\tilde{\tilde\rho}^{\beta}\le
r_k^{\beta} +8\tilde{\tilde\rho}^{\beta}<r_k^{\beta}
+16\tilde\rho_k^{\beta}.
$$
This then implies
$$
\tilde{\tilde\rho}_k^{\beta}=\tilde{\tilde\rho}(r_k^{\beta})=
\tilde\rho(r_k^{\beta}+20\tilde\rho(r_k^{\beta}))
<\tilde\rho(r_k^{\alpha})=\tilde\rho_k^{\alpha}
$$
and hence
$$
\tilde{\tilde\rho}^{\beta}\le 4\tilde\rho^{\alpha}.
$$
Therefore $\hat B_k^{\beta}\subset B(x_k^{\alpha},
36\tilde\rho^{\alpha})$ whenever $\hat B_k^{\alpha}$ and $\hat
B_k^{\beta}$ meet and $k\ge \max\{ k({\alpha}),$ $k({\beta})\}.$

Next we define the balls $B_k^{\alpha}=B(x_k^{\alpha},
5\tilde{\tilde\rho}^{\alpha})$ and $\tilde
B_k^{\alpha}=B(x_k^{\alpha}, \tilde{\tilde\rho}^{\alpha}/2)$. Note
that $\tilde B_k^{\alpha}$ are disjoint since  $\tilde
B_k^{\alpha}\subset B(x_k^{\alpha},
\tilde{\tilde{\rho}}^{\alpha}_k)$. Since $\hat B_k^{\alpha}\subset
B_k^{\alpha}$, the balls $B_k^{\alpha}$ cover $B(p_k, r)$ for
$\alpha\le \lambda(r)$ as before. If $B_k^{\alpha}$ and
$B_k^{\beta}$ meet for $k\ge k({\alpha})$ and $k\ge k({\beta})$
and $r_k^{\beta}\le r_k^{\alpha}$, then by the triangle inequality
we get
$$
r_k^{\alpha}\le r_k^{\beta}
+10\tilde{\tilde\rho}^{\beta}<r_k^{\beta} +20\tilde\rho_k^{\beta},
$$
and hence
$$
\tilde{\tilde\rho}^{\beta}\le 4\tilde\rho^{\alpha}
$$
again. Similarly,
$$
\tilde\rho_k^{\beta}=\tilde\rho(r_k^{\beta})=
\rho(r_k^{\beta}+20\rho(r_k^{\beta}))
<\rho(r_k^{\alpha})=\rho_k^{\alpha}.
$$
This makes
$$
\tilde\rho^{\beta}\le 4\rho^{\alpha}.
$$
Now any point in $B_k^{\beta}$ has distance at most
$$
5\tilde{\tilde\rho}^{\alpha}+5\tilde{\tilde\rho}^{\beta}
+5\tilde{\tilde\rho}^{\beta}\le 45\tilde\rho^{\alpha}
$$
from $x_k^{\alpha}$, so $B_k^{\beta}\subset B(x_k^{\alpha},
45\tilde\rho^{\alpha})$.  Likewise, whenever $B_k^{\alpha}$ and
$B_k^{\beta}$ meet for $k\ge k({\alpha})$ and $k\ge k({\beta})$,
any point in the larger ball $B(x_k^{\beta},
45\tilde\rho^{\beta})$ has distance at most
$$
5\tilde{\tilde\rho}^{\alpha}+5\tilde{\tilde\rho}^{\beta}
+45\tilde\rho^{\beta}\le 205\rho^{\alpha}
$$
from $x_k^{\alpha}$ and hence $B(x_k^{\beta},
45\tilde\rho^{\beta})\subset B(x_k^{\alpha}, 205\rho^{\alpha})$.
Now we define ${\bar B}_k^{\alpha}=B(x_k^{\alpha},
45\tilde\rho^{\alpha})$ and ${\bar{\bar
B}}_k^{\alpha}=B(x_k^{\alpha}, 205\rho^{\alpha})$. Then the above
discussion says that whenever $B_k^{\alpha}$ and $B_k^{\beta}$
meet for $k\ge k({\alpha})$ and $k\ge k({\beta})$, we have \be
B_k^{\beta}\subset {\bar B}_k^{\alpha} \quad \mbox{and}\quad {\bar
B}_k^{\beta}\subset {\bar{\bar
B}}_k^{\alpha}. %\eqno (4.1.6)
\ee Note that ${\bar{\bar B}}_k^{\alpha}$ is still a nice embedded
ball since, by (4.1.3), $205\rho^{\alpha}\le
410\rho_k^{\alpha}<inj(M_k, x_k^{\alpha}). $

We claim there exist positive numbers $N(r)$ and $k''(r)$ such
that for any given $\alpha$ with $r^{\alpha}<r$, as $k \geq
k''(r)$, there holds \be \# \{\beta \ |\ B_k^{\alpha} \cap
B_k^{\beta} \neq \phi \}
\leq N(r). %\eqno (4.1.7)
\ee

Indeed, if $B_k^{\alpha}$ meets $B_k^{\beta}$ then there is a
positive $k''(\alpha)$ such that as $k \geq k''(\alpha)$,
\begin{align*}
r_k^\beta
&  \leq r_k^{\alpha} + 10\tilde{\tilde{\rho}}_k^{\alpha }\\
&  \leq r + 20\rho(r)\\
&  \leq r + \frac{1}{5}(s_{\infty}-r),
\end{align*}
where we used (4.1.2) in the third inequality. Set
$$
k''(r)=\max\{ k''(\alpha), k'(r)\ |\ \alpha \leq \lambda(r)\}
$$ and
$$
N(r)= \lambda\(r + \frac{1}{5}(s_{\infty}-r)\).
$$
Then by combining with (4.1.4), these give the desired estimate
(4.1.7)

Next we observe that by passing to another subsequence we can
guarantee that for any pair $\alpha$ and $\beta$ we can find a
number $k(\alpha,\beta)$ such that if $k\ge k(\alpha,\beta)$ then
either $B_k^{\alpha}$ always meets $B_k^{\beta}$ or it never does.

Hence by setting
\begin{multline*}
\bar{k}(r) = \max \bigg\{ k(\alpha,\beta), k(\alpha), k(\beta),
k''(r) \ |\
\alpha \leq\lambda(r)\; \mbox{ and}\\
\beta \leq \lambda\(r +\frac{1}{5}(s_{\infty}-r)\) \bigg\},
\end{multline*}
we have shown the following results: for every $r < s_{\infty}$,
if $k\ge \bar{k}(r)$, we have
\begin{itemize}
\item[(i)] the ball $B(p_k, r)$ in $M_k$ is covered by the balls
$B_k^{\alpha}$ for $\alpha\le \lambda(r)$, \item[(ii)] whenever
$B_k^{\alpha}$ and $B_k^{\beta}$ meet for $\alpha \leq
\lambda(r)$, we have
$$
B_k^{\beta}\subset {\bar B}_k^{\alpha} \quad \mbox{and}\quad {\bar
B}_k^{\beta}\subset {\bar{\bar B}}_k^{\alpha},
$$
\item[(iii)] for each $\alpha \leq \lambda(r)$, there no more than
$N(r)$ balls ever meet $B_k^{\alpha}$, and \item[(iv)] for any
$\alpha \le \lambda(r)$ and any $\beta$, either $B_k^{\alpha}$
meets $B_k^{\beta}$ for all $k\ge \bar{k}(r)$ or none for all
$k\ge \bar{k}(r)$.
\end{itemize}

Now we let $\hat E^{\alpha}, E^{\alpha}$, $\bar E^{\alpha}$, and
${\bar{\bar E}}^{\alpha}$ be the balls of radii
$4\tilde{\tilde\rho}^{\alpha}, 5\tilde{\tilde\rho}^{\alpha},
45\tilde\rho^{\alpha}$, and $205\rho^{\alpha}$ around the origin
in Euclidean space $\mathbb R^n$. At each point $x_k^{\alpha}\in
B(p_k, s_k)$ we define coordinate charts
$H_k^{\alpha}:E^{\alpha}\to B^{\alpha}_k$ as the composition of a
linear isometry of $\mathbb R^n$ to the tangent space
$T_{x_k^{\alpha}}M_k$ with the exponential map
$\exp_{x_k^{\alpha}}$ at $x_k^{\alpha}$. We also get maps $\bar
H_k^{\alpha}:\bar E^{\alpha}\to \bar B^{\alpha}_k$ and ${\bar{\bar
H}}_k^{\alpha}:{\bar{\bar E}}^{\alpha}\to {\bar{\bar
B}}^{\alpha}_k$ in the same way. Note that (4.1.3) implies that
these maps are all well defined. We denote by $g_k^{\alpha}$ (and
$\bar g_k^{\alpha}$ and ${\bar{\bar g}}_k^{\alpha}$) the
pull-backs of the metric $g_k$ by $H_k^{\alpha}$ (and $\bar
H_k^{\alpha}$ and ${\bar{\bar H}}_k^{\alpha}$). We also consider
the coordinate transition functions $J_k^{\alpha\beta}:
E^{\beta}\to \bar E^{\alpha}$ and $\bar J_k^{\alpha\beta}: \bar
E^{\beta}\to {\bar{\bar E}}^{\alpha}$ defined by
$$
J_k^{\alpha\beta}=(\bar H_k^{\alpha})^{-1}H_k^{\beta} \quad
{\text{and}} \quad \bar J_k^{\alpha\beta}=({\bar{\bar
H}}_k^{\alpha})^{-1}\bar H_k^{\beta}.
$$
Clearly $\bar J_k^{\alpha\beta} J_k^{\beta\alpha}=I$. Moreover
$J_k^{\alpha\beta}$ is an isometry from $g_k^{\beta}$ to $\bar
g_k^{\alpha}$ and $\bar J_k^{\alpha\beta}$ from $\bar g_k^{\beta}$
to ${\bar{\bar g}}_k^{\alpha}$.

Now for each fixed $\alpha$, the metrics $g_k^{\alpha}$ are in
geodesic coordinates and have their curvatures and their covariant
derivatives uniformly bounded.

% CZ-EXTRA source-line:7290 source-id:theorem-4.1.2-claim-1
\end{proof}
\begin{proposition}[local metric subsequence convergence]
  \label{prop:cz-local-metric-subsequence-convergence}
  \dcref{Claim 1}
  \group{Cao--Zhu Cheeger Type Compactness}
  \source{cao-zhu}
By passing to another subsequence we
can guarantee that for each $\alpha$ (and indeed all $\alpha$ by
diagonalization) the metrics $g_k^{\alpha}$ (or $\bar
g_k^{\alpha}$ or ${\bar{\bar g}}_k^{\alpha}$) converge uniformly
with their derivatives to a smooth metric $g^{\alpha}$ (or $\bar
g^{\alpha}$ or ${\bar{\bar g}}^{\alpha}$) on $E^{\alpha}$ (or
$\bar E^{\alpha}$ or ${\bar{\bar E}}^{\alpha}$) which is also in
geodesic coordinates.
\end{proposition}
\begin{proof}
\ It suffices to show the
following general result:

\medskip
{\it There exists a constant $c>0$ depending only on the
dimension, and constants $C_q$ depending only on the dimension and
$q$ and bounds $B_j$ on the curvature and its derivatives for $j
\leq q$ where $|D^jRm| \leq B_j$, so that for any metric $g_{k
\ell}$ in geodesic coordinates in the ball $|x| \leq r \leq c/
\sqrt{B_0}$, we have
$$
\frac{1}{2} I_{k \ell} \leq g_{k \ell} \leq 2I_{k \ell}
$$
and
$$
\Big| \frac{\partial}{\partial x^{j_1}} \cdots
\frac{\partial}{\partial x^{j_q}} g_{k \ell} \Big| \leq C_q,
$$
where $I_{k\ell}$ is the Euclidean metric.}

\medskip
Suppose we are given a metric $g_{ij}(x)dx^idx^j$ in geodesic
coordinates in the ball $|x| \leq r \leq c/ \sqrt{B_0}$ as in
Claim 1. Then by definition every line through the origin is a
geodesic (parametrized proportional to arc length) and
$g_{ij}=I_{ij}$ at the origin. Also, the Gauss Lemma says that the
metric $g_{ij}$ is in geodesic coordinates if and only if
$g_{ij}x^i=I_{ij}x^i$. Note in particular that in geodesic
coordinates
$$
|x|^2=g_{ij}x^ix^j=I_{ij}x^ix^j
$$
is unambiguously defined. Also, in geodesic coordinates we have
$\Gamma^k_{ij}(0)=0$, and all the first derivatives for $g_{jk}$
vanish at the origin.

Introduce the symmetric tensor
$$
A_{ij} = \frac{1}{2} x^k \frac{\partial}{\partial x^k} g_{ij}.
$$
Since we have $g_{jk}x^k=I_{jk}x^k$, we get
$$
x^k \frac{\partial}{\partial x^i}g_{jk} = I_{ij}-g_{ij}=x^k
\frac{\partial}{\partial x^j}g_{ik}
$$
and hence from the formula for $\Gamma^i_{jk}$
$$
x^j \Gamma^i_{jk}=g^{i \ell}A_{k \ell}.
$$
Hence $A_{k \ell}x^k=0$. Let $D_i$ be the covariant derivative
with respect to the metric $g_{ij}$. Then
$$
D_ix^k=I^k_i+ \Gamma^k_{ij}x^j=I^k_i+g^{k \ell}A_{i \ell}.
$$
Introduce the potential function
$$
P=|x|^2/2 = \frac{1}{2}g_{ij}x^ix^j.
$$
We can use the formulas above to compute
$$
D_iP=g_{ij}x^j.
$$
Also we get
$$
D_iD_jP=g_{ij}+A_{ij}.
$$
The defining equation for $P$ gives
$$
g^{ij}D_iPD_jP=2P.
$$
If we take the covariant derivative of this equation we get
$$
g^{k \ell}D_jD_kPD_{\ell}P=D_jP
$$
which is equivalent to $A_{jk}x^k=0$. But if we take the covariant
derivative again we get
$$
g^{k \ell}D_iD_jD_kPD_\ell P + g^{k \ell}D_jD_kPD_iD_\ell P = D_i
D_j P.
$$
Now switching derivatives
$$
D_iD_jD_kP=D_iD_kD_jP=D_kD_iD_jP+R_{ikj \ell}g^{\ell m}D_mP
$$
and if we use this and $D_iD_jP=g_{ij}+A_{ij}$ and $g^{k
\ell}D_{\ell}P=x^k$ we find that
$$
x^kD_kA_{ij}+A_{ij}+g^{k \ell} A_{ik}A_{j \ell} + R_{ikj \ell}x^k
x^\ell=0.
$$
{}From our assumed curvature bounds we can take $|R_{ijk \ell}|
\leq B_0$. Then we get the following estimate:
$$
|x^kD_kA_{ij}+A_{ij}| \leq C|A_{ij}|^2+CB_0 r^2
$$
on the ball $|x| \leq r$ for some constant $C$ depending only on
the dimension.

We now show how to use the maximum principle on such equations.
First of all, by a maximum principle argument, it is easy to show
that if $f$ is a function on a ball $|x| \leq r$ and $\lambda>0$
is a constant, then
$$
\lambda \operatorname{sup}|f| \leq \operatorname{sup} \Big| x^k
\frac{\partial f}{\partial x^k} + \lambda f \Big|.
$$
For any tensor $T=\{T_{i \cdots j}\}$ and any constant $\lambda
>0$, setting $f=|T|^2$ in the above inequality, we have \be
\lambda \sup|T| \leq \sup|x^kD_kT+
\lambda T|. %\eqno (4.1.8)
\ee Applying this to the tensor $A_{ij}$ we get
\[
\sup_{|x| \leq r} |A_{ij}| \leq C \sup_{|x| \leq r}
|A_{ij}|^2+CB_0r^2
\]
for some constant depending only on the dimension.

It is fairly elementary to see that there exist constants $c>0$
and $C_0 < \infty$ such that if the metric $g_{ij}$ is in geodesic
coordinates with $|R_{ijk \ell}| \leq B_0$ in the ball of radius
$r \leq c /\sqrt{B_0}$ then
$$
|A_{ij}| \leq C_0 B_0r^2.
$$
Indeed, since the derivatives of $g_{ij}$ vanish at the origin, so
does $A_{ij}$. Hence the estimate holds near the origin. But the
inequality
$$
\sup_{|x| \leq r} |A_{ij}| \leq C \sup_{|x| \leq r}
|A_{ij}|^2+CB_0r^2
$$
says that $|A_{ij}|$ avoids an interval when $c$ is chosen small.
In fact the inequality
$$
X \leq CX^2+D
$$
is equivalent to
$$
|2CX-1| \geq \sqrt{1-4CD}
$$
which makes $X$ avoid an interval if $4CD<1$. (Hence in our case
we need to choose $c$ with $4C^2c^2<1.)$ Then if $X$ is on the
side containing $0$ we get
$$
X \leq \frac{1- \sqrt{1 - 4CD}}{2C} \leq 2D.
$$
This gives $|A_{ij}| \leq C_0B_0r^2$ with $C_0=2C$.

\smallskip
We can also derive bounds on all the covariant derivatives of $P$
in terms of bounds on the covariant derivatives of the curvature.
To simplify the notation, we let
$$
D^qP= \{D_{j_1}D_{j_2} \cdots D_{j_q}P \}
$$
denote the $q^{\operatorname{th}}$ covariant derivative, and in
estimating $D^qP$ we will lump all the lower order terms into a
general slush term $\Phi^q$ which will be a polynomial in
$D^1P,D^2P, \ldots, D^{q-1}P$ and $Rm, D^1Rm, \ldots, D^{q-2}Rm$.
We already have estimates on a ball of radius $r$
$$
P \leq r^2/2
$$
$$
|D^1P|\leq r
$$
$$
|A_{ij}| \leq C_0 B_0r^2
$$
and since $D_iD_jP=g_{ij}+A_{ij}$ and $r \leq c /\sqrt{B_0}$ if we
choose $c$ small we can make
$$
|A_{ij}|\leq 1/2,
$$
and we get
$$
|D^2P| \leq C_2
$$
for some constant $C_2$ depending only on the dimension.

\smallskip
Start with the equation $g^{ij}D_iPD_jP=2P$ and apply repeated
covariant derivatives. Observe that we get an equation which
starts out
$$
g^{ij}D_iPD^qD_jP+ \cdots =0
$$
where the omitted terms only contain derivatives $D^qP$ and lower.
If we switch two derivatives in a term $D^{q+1}P$ or lower, we get
a term which is a product of a covariant derivative of $Rm$ of
order at most $q-2$ (since the two closest to $P$ commute) and a
covariant derivative of $P$ of order at most $q-1$; such a term
can be lumped in with the slush term $\Phi^q$. Therefore up to
terms in $\Phi^q$ we can regard the derivatives as commuting. Then
paying attention to the derivatives in $D^1P$ we get an equation

\smallskip\noindent
$g^{ij}D_iPD_jD_{k_1} \cdots D_{k_q}P+g^{ij}D_iD_{k_1}PD_jD_{k_2}
\cdots D_{k_q}P$

\smallskip\noindent
$+g^{ij}D_iD_{k_2}PD_jD_{k_1}D_{k_3} \cdots D_{k_q}P + \cdots +
g^{ij}D_iD_{k_q}PD_jD_{k_1} \cdots D_{k_{q-1}}P$

\smallskip\noindent
$ =D_{k_1} \cdots D_{k_q}P+ \Phi^q$.

\medskip
Recalling that $D_iD_jP=g_{ij}+A_{ij}$ we can rewrite this as
\begin{align*}
\Phi^q &  = g^{ij}D_iPD_jD_{k_1} \cdots D_{k_q}+(q-1)D_{k_1}
\cdots
D_{k_q}P \\
&\quad  + g^{ij}A_{ik_1}D_jD_{k_2} \cdots D_{k_q}P+ \cdots +
g^{ij}A_{ik_q}D_jD_{k_1} \cdots D_{k_{q-1}}P.
\end{align*}
Estimating the product of tensors in the usual way gives
$$
|x^iD_iD^qP+(q-1)D^qP| \leq q|A||D^qP|+|\Phi^q|.
$$
Applying the inequality $\lambda \sup|T| \leq \sup|x^kD_kT+
\lambda T|$ with $T=D^qP$ gives
$$
(q-1) \sup|D^qP| \leq\sup(q|A||D^qP|+| \Phi^q|).
$$
Now we can make $|A| \leq 1/2$ by making $r \leq c/ \sqrt{B_0}$
with $c$ small; it is important here that $c$ is independent of
$q!$ Then we get
$$
(q-2) \operatorname{sup}|D^qP| \leq 2 \operatorname{sup} | \Phi^q|
$$
which is a good estimate for $q \geq 3$. The term $\Phi^q$ is
estimated inductively from the terms $D^{q-1}P$ and $D^{q-2}Rm$
and lower. This proves that there exist constants $C_q$ for $q
\geq 3$ depending only on $q$ and the dimension and on $|D^jRm|$
for $j \leq q-2$ such that
$$
|D^qP| \leq C_q
$$
on the ball $r \leq c/ \sqrt{B_0}.$

Now we turn our attention to estimating the Euclidean metric
$I_{jk}$ and its covariant derivatives with respect to $g_{jk}$.
We will need the following elementary fact: suppose that $f$ is a
function on a ball $|x| \leq r$ with $f(0)=0$ and
$$
\Big| x^i \frac{\partial f}{\partial x^i} \Big| \leq C |x|^2
$$
for some constant $C$. Then \be
|f| \leq C|x|^2  %\eqno (4.1.9)
\ee for the same constant $C$. As a consequence, if $T = \{T_{j
\cdots k }\}$ is a tensor which vanishes at the origin and if
$$
|x^iD_iT| \leq C|x|^2
$$
on a ball $|x| \leq r$ then $|T| \leq  C|x|^2$ with the same
constant $C$. (Simply apply the inequality (4.1.9) to the function
$f=|T|$. In case this is not smooth, we can use
$f=\sqrt{|T|^2+\epsilon^2}-\epsilon$ and then let $\epsilon\to
0$.)

Our application will be to the tensor $I_{jk}$ which gives the
Euclidean metric as a tensor in geodesic coordinates. We have
$$
D_iI_{jk} = - \Gamma^p_{ij}I_{pk} - \Gamma^p_{ik}I_{pj}
$$
and since
$$
x^i \Gamma^p_{ij} = g^{pq}A_{jq}
$$
we get the equation
$$
x^iD_iI_{jk} = -g^{pq}A_{jp}I_{kq}-g^{pq}A_{kp}I_{jq}
$$.

\noindent We already have $|A_{jk}| \leq C_0B_0|x|^2$ for $|x|
\leq r \leq c/ \sqrt{B_0}$. The tensor $I_{jk}$ doesn't vanish at
the origin, but the tensor
$$
h_{jk}=I_{jk}-g_{jk}
$$
does. We can then use
$$
x^iD_ih_{jk}=-g^{pq}A_{jp}h_{kq}-g^{pq}A_{kq}h_{jq}-2A_{jk}.
$$
Suppose $M(s) = \sup_{|x| \leq s}|h_{jk}|$. Then
$$
|x^iD_ih_{jk}| \leq 2[1+M(s)]C_0B_0|x|^2
$$
and we get
$$
|h_{jk}| \leq 2 [1+M(s)]C_0B_0|x|^2
$$
on $|x| \leq s$. This makes
$$
M(s) \leq 2[1+M(s)]C_0B_0s^2.
$$
Then for $s \leq r \leq c/ \sqrt{B_0}$ with $c$ small compared to
$C_0$ we get $2C_0B_0s^2 \leq 1/2$ and $M(s) \leq 4C_0B_0s^2$.
Thus
$$
|I_{jk}-g_{jk}|=|h_{jk}| \leq 4C_0B_0|x|^2
$$
for $|x| \leq r\leq c/ \sqrt{B_0}$, and hence for $c$ small enough
$$
\frac{1}{2} g_{jk} \leq I_{jk} \leq 2g_{jk}.
$$
Thus the metrics are comparable. Note that this estimate only
needs $r$ small compared to $B_0$ and does not need any bounds on
the derivatives of the curvature.

Now to obtain bounds on the covariant derivative of the Eucliden
metric $I_{k \ell}$ with respect to the Riemannian metric $g_{k
\ell}$ we want to start with the equation
$$
x^iD_iI_{k \ell}+g^{mn}A_{km}I_{\ell n}+g^{mn}A_{\ell m}I_{kn}=0
$$
and apply $q$ covariant derivatives $D_{j_1} \cdots D_{j_q}$. Each
time we do this we must interchange $D_j$ and $x^iD_i$, and since
this produces a term which helps we should look at it closely. If
we write $R_{ji}=[D_j,D_i]$ for the commutator, this operator on
tensors involves the curvature but no derivatives. Since
$$
D_jx^i=I^i_j+g^{im}A_{jm}
$$
we can compute
$$
[D_j, x^iD_i]=D_j+g^{im}A_{jm}D_i+x^iR_{ji}
$$
and the term $D_j$ in the commutator helps, while $A_{jm}$ can be
kept small and $R_{ji}$ is zero order. It follows that we get an
equation of the form
\begin{align*}
0 &  = x^iD_iD_{j_1} \cdots D_{j_q}I_{k \ell} + qD_{j_1} \cdots
D_{j_q}I_{k \ell} \\
&\quad  + \sum^q_{h=1} g^{im}A_{j_hm}D_{j_1} \cdots
D_{j_{h-1}}D_iD_{j_{h+1}} \cdots D_{j_q}I_{k \ell} \\
&\quad  + g^{mn}A_{km}D_{j_1} \cdots D_{j_q}I_{\ell n} + g^{mn}
A_{\ell m} D_{j_1} \cdots D_{j_q}I_{kn}+ \Psi^q,
\end{align*}
where the slush term $\Psi^q$ is a polynomial in derivatives of
$I_{k \ell}$ of degree no more than $q-1$ and derivatives of $P$
of degree no more than $q+2$ (remember $x^i = g^{ij}D_jP$ and
$A_{ij} = D_iD_jP-g_{ij})$ and derivatives of the curvature $Rm$
of degree no more than $q-1$. We now estimate
$$
D^qI_{k \ell} = \{D_{j_1} \cdots D_{j_q}I_{k \ell}\}
$$
by induction on $q$ using (4.1.8) with $\lambda = q$. Noticing a
total of $q+2$ terms contracting $A_{ij}$ with a derivative of
$I_{k \ell}$ of degree $q$, we get the estimate
$$
q \sup |D^qI_{k \ell}| \leq (q+2) \operatorname{sup}|A|
\operatorname{sup} |D^q I_{k \ell} | + \operatorname{sup} |
\Psi^q|.
$$

\noindent and everything works. This proves that there exists a
constant $c>0$ depending only on the dimension, and constants
$C_q$ depending only on the dimension and $q$ and bounds $B_j$ on
the curvature and its derivatives for $j \leq q$ where $|D^jRm|
\leq B_j$, so that for any metric $g_{k \ell}$ in geodesic
coordinates in the ball $|x| \leq r \leq c/ \sqrt{B_0}$ the
Euclidean metric $I_{k \ell}$ satisfies
$$
\frac{1}{2} g_{k \ell}\leq I_{k \ell} \leq 2g_{k \ell}
$$
and the covariant derivatives of $I_{k \ell}$ with respect to
$g_{k \ell}$ satisfy
$$
|D_{j_1} \cdots D_{j_q}I_{k \ell}| \leq C_q.
$$

The difference between a covariant derivative and an ordinary
derivative is given by the connection
$$
-\Gamma^p_{ij}I_{pk} - \Gamma^p_{ik} I_{pj}
$$
to get
$$
\Gamma^k_{ij} = \frac{1}{2} I^{k \ell}(D_\ell I_{ij} - D_iI_{j
\ell} - D_j I_{i \ell}).
$$
This gives us bounds on $\Gamma^k_{ij}$. We then obtain bounds on
the first derivatives of $g_{ij}$ from
$$
\frac{\partial}{\partial x^i} g_{jk} = g_{k \ell} \Gamma^\ell_{ij}
+ g_{j \ell} \Gamma^\ell_{ik}.
$$
Always proceeding inductively on the order of the derivative, we
now get bounds on covariant derivatives of $\Gamma^k_{ij}$ from
the covariant derivatives of $I_{pk}$ and bounds of the ordinary
derivatives of $\Gamma^k_{ij}$ by relating the to the covariant
derivatives using the $\Gamma^k_{ij}$, and bounds on the ordinary
derivatives of the $g_{jk}$ from bounds on the ordinary
derivatives of the $\Gamma^\ell_{ij}$. Consequently, we have
estimates
$$
\frac{1}{2} I_{k \ell} \leq g_{k \ell} \leq 2I_{k \ell}
$$
and
$$
\Big| \frac{\partial}{\partial x^{j_1}} \cdots
\frac{\partial}{\partial x^{j_q}} g_{k \ell} \Big| \leq
\tilde{C}_q
$$
for similar constants $\tilde{C}_q$.

Therefore we have finished the proof of Claim 1.
\end{proof}
\begin{proof}
  \proves{thm:cz-cheeger-type-pointed-compactness}

Look now at any pair $\alpha, \beta$ for which the balls
$B_k^{\alpha}$ and $B_k^{\beta}$ always meet for large $k$, and
thus the maps $J_k^{\alpha\beta}$ (and $\bar J_k^{\alpha\beta}$
and $J_k^{\beta\alpha}$ and $\bar J_k^{\beta \alpha}$) are always
defined for large $k$.

% CZ-EXTRA source-line:7305 source-id:theorem-4.1.2-claim-2
\end{proof}
\begin{proposition}[transition isometry subsequence convergence]
  \label{prop:cz-transition-isometry-subsequence-convergence}
  \dcref{Claim 2}
  \group{Cao--Zhu Cheeger Type Compactness}
  \source{cao-zhu}
The isometries $J_k^{\alpha\beta}$
(and $\bar J_k^{\alpha\beta}$ and $J_k^{\beta\alpha}$ and $\bar
J_k^{\beta\alpha}$) always have a convergent subsequence. \vskip
0.2cm
\end{proposition}
\begin{proof}
\  We need to show how to
estimate the derivatives of an isometry. We will prove that if
$y=F(x)$ is an isometry from a ball in Euclidean space with a
metric $g_{ij}dx^idx^j$ to a ball in Euclidean space with a metric
$h_{kl}dy^kdy^l$. Then we can bound all of the derivatives of $y$
with respect to $x$ in terms of bounds on $g_{ij}$ and its
derivatives with respect to $x$ and bound on $h_{kl}$ and its
derivatives with respect to $y$. This would imply Claim 2.

Since $y=F(x)$ is an isometry we have the equation
$$
h_{pq} \frac{\partial y^p}{\partial x^j} \frac{\partial
y^q}{\partial x^k} = g_{jk}.
$$

\noindent Using bounds $g_{jk} \leq CI_{jk}$ and $h_{pq} \geq
cI_{pq}$ comparing to the Euclidean metric, we easily get
estimates
$$
\Big| \frac{\partial y^p}{\partial x^j} \Big| \leq C.
$$

\noindent Now if we differentiate the equation with respect to
$x^i$ we get
$$
h_{pq} \frac{\partial^2y^p}{\partial x^i \partial x^j}
\frac{\partial y^q}{\partial x^k} + h_{pq} \frac{\partial
y^p}{\partial x^j} \frac{\partial^2y^q}{\partial x^i \partial x^k}
= \frac{\partial g_{jk}}{\partial x^i} - \frac{\partial
h_{pq}}{\partial y^r} \frac{\partial y^r}{\partial x^i}
\frac{\partial y^p}{\partial x^j} \frac{\partial y^q}{\partial
x^k}.
$$

\noindent Now let
$$
T_{ijk} = h_{pq} \frac{\partial y^p}{\partial x^i}
\frac{\partial^2 y^q}{\partial x^j \partial x^k}
$$

\noindent and let
$$
U_{ijk} = \frac{\partial g_{jk}}{\partial x^i} - \frac{\partial
h_{pq}}{\partial y^r} \frac{\partial y^r}{\partial x^i}
\frac{\partial y^p}{\partial x^j} \frac{\partial y^q}{\partial
x^k}.
$$

\noindent Then the above equation says
$$
T_{kij}+T_{jik}=U_{ijk}.
$$

\noindent Using the obvious symmetries $T_{ijk}=T_{ikj}$ and
$U_{ijk} = U_{ikj}$ we can solve this in the usual way to obtain
$$
T_{ijk} = \frac{1}{2}(U_{jik}+U_{kij} - U_{ijk}).
$$

\noindent We can recover the second derivatives of $y$ with
respect to $x$ from the formula
$$
\frac{\partial^2 y^p}{\partial x^i \partial x^j} = g^{k
\ell}T_{kij} \frac{\partial y^p}{\partial x^\ell}.
$$

\noindent Combining these gives an explicit formula giving
$\partial^2y^p/\partial x^i \partial x^j$ as a function of
$g^{ij}, h_{pq}, \partial g_{jk}/\partial x^i, \partial
h_{pq}/\partial y^r$, and $ \partial y^p/\partial y^i$. This gives
bounds
$$
\Big| \frac{\partial^2 y^p}{\partial y^i \partial y^j}\Big| \leq C
$$

\noindent and bounds on all higher derivatives follow by
differentiating the formula and using induction. This completes
the proof of Claim 2 and hence the proof of Theorem 4.1.2.
\end{proof}
\begin{proof}
  \proves{thm:cz-cheeger-type-pointed-compactness}

So by passing to another subsequence we may assume
$J_k^{\alpha\beta}\to J^{\alpha\beta}$ (and $\bar
J_k^{\alpha\beta}\to \bar J^{\alpha\beta}$ and
$J_k^{\beta\alpha}\to J^{\beta\alpha}$ and $\bar
J_k^{\beta\alpha}\to \bar J^{\beta\alpha}$). The limit maps
$J^{\alpha\beta}: E^{\beta}\to \bar E^{\alpha}$ and $\bar
J^{\alpha\beta}: \bar E^{\beta}\to {\bar{\bar E}}^{\alpha}$ are
isometries in the limit metrics $g^{\beta}$ and $g^{\alpha}$.
Moreover
$$
J^{\alpha\beta}\bar J^{\beta\alpha}=I.
$$
We are now done picking subsequences, except we still owe the
reader the proofs of Claim 1 and Claim 2.

\medskip
{\it Step} 2: Finding local diffeomorphisms which are approximate
isometries.

\smallskip
Take the subsequence $(B(p_k, s_k), g_k, p_k)$ chosen in Step 1
above. We claim that for every $r < s_{\infty}$ and every
$(\epsilon_1,\epsilon_2,\ldots,\epsilon_p)$, and for all $k$ and
$l$ sufficiently large in comparison, we can find a diffeomorphism
$F_{kl}$ of a neighborhood of the ball $B(p_k, r)\subset B(p_k,
s_k)$ into an open set in $B(p_l, s_l)$ which is an
$(\epsilon_1,\epsilon_2,\ldots,\epsilon_p)$ approximate isometry
in the sense that
$$
|^{t}\nabla F_{kl} \nabla F_{kl}-I|<\epsilon_1
$$
and
$$
|\nabla^2F_{kl}|<\epsilon_2,\ldots, |\nabla^pF_{kl}|<\epsilon_p
$$
where $\nabla^pF_{kl}$ is the $p^{th}$ covariant derivative of
$F_{kl}$.

The idea (following Peters \cite{Pe} or Greene and Wu \cite{Gre})
of proving the claim is to define the map $F^\alpha_{k
\ell}=H^{\alpha}_l\circ(H^{\alpha}_k)^{-1}$ (or $\bar{F}^\alpha_{k
\ell}=\bar{H}^{\alpha}_l\circ(\bar{H}^{\alpha}_k)^{-1}$, resp.)
from $B^\alpha_k$ to $B^\alpha_\ell$ (or $\bar{B}^\alpha_k$ to
$\bar{B}^\alpha_\ell$, resp.) for $k$ and $\ell$ large compared to
$\alpha$ so as to be the identity map on ${E}^\alpha$ (or
$\bar{E}^\alpha$, resp.) in the coordinate charts $H^\alpha_k$ and
$H^\alpha_\ell$ (or $\bar{H}^\alpha_k$ and $\bar{H}^\alpha_\ell$,
resp.), and then to define $F_{k \ell}$ on a neighborhood of
$B(p_k,r)$ for $k, \ell \geq \bar{k}(r)$ be averaging the maps
$\bar{F}^\beta_{k \ell}$ for $\beta \leq \
\lambda(r+\frac{1}{5}(s_{\infty}-r))$. To describe the averaging
process on $B^\alpha_k$ with $\alpha \leq \lambda(r)$ we only need
to consider those $B^\beta_k$ which meet $B^\alpha_k$; there are
never more than $N(r)$ of them and each $\beta \leq
\lambda(r+\frac{1}{5}(s_{\infty}-r))$, and they are the same for
$k$ and $\ell$ when $k, \ell \geq \bar{k}(r)$. The averaging
process is defined by taking $F_{k \ell}(x)$ to be the center of
mass of the $\bar{F}^\beta_{k \ell}(x)$ for $x \in B^\alpha_k$
averaging over those $\beta$ where $B^\beta_k$ meets $B^\alpha_k$
using weights $\mu^\beta_k(x)$ defined by a partition of unity.
The center of mass of the points $y^\beta=F^\beta_{k \ell}(x)$
with weights $\mu^\beta$ is defined to be the point $y$ such that
$$
\exp_yV^\beta = y^\beta \quad \text{and} \quad \sum \mu^\beta
V^\beta=0.
$$
When the points $y^\beta$ are all close and the weights
$\mu^\beta$ satisfy $0 \leq \mu^\beta \leq 1$ then there will be a
unique solution $y$ close to $y^\beta$ which depends smoothly on
the $y^\beta$ and the $\mu^\beta$ (see \cite{Gre} for the
details). The point $y$ is found by the inverse function theorem,
which also provides bounds on all the derivatives of $y$ as a
function of the $y^\beta$ and the $\mu^\beta$.

Since $B^\alpha_k \subseteq \bar{B}^\beta_k$ and
$\bar{B}^\beta_\ell \subseteq \bar{\bar{B}}^{\alpha}_\ell$, the
map $\bar{F}^\beta_{k
\ell}=\bar{H}^{\beta}_l\circ(\bar{H}^{\beta}_k)^{-1}$ can be
represented in local coordinates by the map
$$
P^{\alpha \beta}_{k \ell}: E^\alpha \rightarrow
\bar{\bar{E}}^\alpha
$$
defined by
$$
P^{\alpha \beta}_{k \ell} = \bar{J}^{\alpha \beta}_\ell \circ
J^{\beta \alpha}_k.
$$
Since $J^{\beta \alpha}_k \rightarrow J^{\beta \alpha}$ as $k
\rightarrow \infty$ and $\bar{J}^{\alpha \beta}_\ell \rightarrow
\bar{J}^{\alpha \beta}$ as $\ell \rightarrow \infty$ and
$\bar{J}^{\alpha \beta} \circ J^{\beta \alpha} = I$, we see that
the maps $P^{\alpha \beta}_{k \ell} \rightarrow I$ as $k, \ell
\rightarrow \infty$ for each choice of $\alpha$ and $\beta$. The
weights $\mu^\beta_k$ are defined in the following way. We pick
for each $\beta$ a smooth function $\psi^\beta$ which equals 1 on
$\hat E^\beta$ and equals 0 outside $E^\beta$. We then transfer
$\psi^\beta$  to a function $\psi^\beta_k$ on $M_k$ by the
coordinate map $\bar{\bar{H}}^\beta_k$ (i.e. $\psi^\beta_k =
\psi^\beta\circ(\bar{\bar{H}}^\beta_k)^{-1}$). Then let
$$
\mu^\beta_k = \psi^\beta_k \Big/ \sum_\gamma \psi^\gamma_k
$$
as usual. In the coordinate chart $E^\alpha$ the function
$\psi^\beta_k$ looks like the composition of $J^{\beta \alpha}_k$
with $\psi^\beta$. Call this function
$$
\psi^{\alpha \beta}_k = \psi^\beta \circ J^{\beta \alpha}_k.
$$
Then as $k \rightarrow \infty, \psi^{\alpha \beta}_k \rightarrow
\psi^{\alpha \beta}$ where
$$
\psi^{\alpha \beta} = \psi^\beta \circ J^{\beta \alpha}.
$$
In the coordinate chart $E^\alpha$ the function $\mu^\beta_k$
looks like
$$
\mu_k^{\alpha \beta} = \psi_k^{\alpha \beta} \Big/ \sum_\gamma
\psi_k^{\alpha \gamma}
$$
and $\mu^{\alpha \beta}_k \rightarrow \mu^{\alpha \beta}$ as $k
\rightarrow \infty$ where
$$
\mu^{\alpha \beta} = \psi^{\alpha \beta} \Big/ \sum_\gamma
\psi^{\alpha \gamma}.
$$
Since the sets $\hat{B}^\alpha_k$ cover $B(p_k,r)$, it follows
that $\sum_\gamma \psi^\gamma_k \geq 1$ on this set and by
combining with (4.1.5) and (4.1.7) there is no problem bounding
all these functions and their derivatives. There is a small
problem in that we want to guarantee that the averaged map still
takes $p_k$ to $p_\ell$. This is true at least for the map $F^0_{k
\ell}$. Therefore it will suffice to guarantee that
$\mu^\alpha_k=0$ in a neighborhood of $p_k$ if $\alpha \neq 0$.
This happens if the same is true for $\psi^\alpha_k$. If not, we
can always replace $\psi^\alpha_k$ by $\tilde{\psi}^\alpha_k=(1-
\psi^0_k)\psi^\alpha_k$ which still leaves $\tilde{\psi}^\alpha_k
\geq \frac{1}{2} \psi^\alpha_k$ or $\psi^0_k \geq \frac{1}{2}$
everywhere, and this is sufficient to make $\sum_\gamma
\tilde{\psi}^\gamma_k \geq \frac{1}{2}$ everywhere.

Now in the local coordinate $E^\alpha$ we are averaging maps
$P^{\alpha \beta}_{k \ell}$ which converge to the identity with
respect to weights $\mu^{\alpha \beta}_k$ which converge. It
follows that the averaged map converges to the identity in these
coordinates. Thus $F_{k \ell}$ can be made to be an $(\epsilon_1,
\epsilon_2, \ldots, \epsilon_p)$ approximate isometry on
$B(p_k,r)$ when $k$ and $\ell$ are suitably large. At least the
estimates
$$
|^t\nabla F_{k \ell} \cdot \nabla F_{k\ell}-I| < \epsilon_1
$$
and $|\nabla^2F_{k \ell}| < \epsilon_2, \ldots, |\nabla^pF_{k
\ell}| < \epsilon_p$ on $B(p_k,r)$ follow from the local
coordinates. We still need to check that $F_{k \ell}$ is a
diffeomorphism on a neighborhood of $B(p_k,r)$.

This, however, follows quickly enough from the fact that we also
get a map $F_{\ell k}$ on a slightly larger ball $B(p_\ell,
r^\prime)$ which contains the image of $F_{k \ell}$ on $B(p_k,r)$
if we take $r^\prime=(1+ \epsilon_1)r$, and $F_{\ell k}$ also
satisfies the above estimates. Also $F_{k \ell}$ and $F_{\ell k}$
fix the markings, so the composition $F_{\ell k} \circ F_{k \ell}$
satisfies the same sort of estimates and fixes the origin $p_k$.

Since the maps $P^{\alpha \beta}_{k \ell}$  and $P^{\alpha
\beta}_{\ell k}$ converge to the identity as $k, \ell$ tend to
infinity, $F_{\ell k} \circ F_{k \ell}$ must be very close to the
identity on $B(p_k,r)$. It follows that $F_{k \ell}$ is
invertible.  This finishes the proof of the claim and the Step 2.

\medskip
{\it Step} 3: Constructing the limit geodesic ball $(B_{\infty},
g_{\infty}, p_{\infty})$.

\smallskip
We now know the geodesic balls $(B(p_k,s_k),g_k,p_k)$ are nearly
isometric for large $k$. We are now going to construct the limit
$B_{\infty}$.  For a sequence of positive numbers $r_j \nearrow
s_{\infty}$ with each $r_j < s_j$, we choose the numbers
$(\epsilon_1(r_j), \ldots, \epsilon_j(r_j))$ so small that when we
choose $k(r_j)$ large in comparison and find the maps
$F_{k(r_j),k(r_{j+1})}$ constructed above on neighborhoods of
$B(p_{k(r_j)},r_j)$, in $M_{k(r_j)}$ into $M_{k(r_{j+1})}$ the
image always lies in $B(p_{k(r_{j+1})}, r_{j+1})$ and the
composition of $F_{k(r_j),k(r_{j+1})}$ with
$F_{k(r_{j+1}),k(r_{j+2})}$ and $\cdots$ and
$F_{k(r_{s-1}),k(r_s)}$ for any $s> j$ is still an $(\eta_1(r_j),
\ldots, \eta_j(r_j))$ isometry for any choice of $\eta_i(r_j)$,
say $\eta_i(r_j)=1/j$ for $1 \leq i \leq j$.  Now we simplify the
notation by writing $M_j$ in place of $M_{k(r_j)}$ and $F_j$ in
place of $F_{k(r_j),k(r_{j+1})}$. Then
$$
F_j:B(p_j,r_j) \rightarrow B(p_{j+1}, r_{j+1})
$$
is a diffeomorphism map from $B(p_j,r_j)$ into $B(p_{j+1},
r_{j+1})$, and the composition
$$
F_{s-1} \circ \cdots \circ F_j: B(p_j,j) \rightarrow B(p_s,s)
$$
is always an $(\eta_1(r_j), \ldots, \eta_j(r_j))$ approximate
isometry.

We now construct the limit $B_{\infty}$ as a topological space by
identifying the balls $B(p_j,r_j)$ with each other using the
homeomorphisms $F_j$. Given any two points $x$ and $y$ in
$B_{\infty}$, we have $x \in B(p_j,r_j)$ and $y \in B(p_s,r_s)$
for some $j$ and $s$. If $j \leq s$ then $x \in B(p_s,r_s)$ also,
by identification. A set in $B_{\infty}$ is open if and only if it
intersects each $B(p_j,r_j)$ in an open set. Then choosing
disjoint neighborhoods of $x$ and $y$ in $B(p_s,r_s)$ gives
disjoint neighborhoods of $x$ and $y$ in $B_{\infty}$. Thus
$B_{\infty}$ is a Hausdorff space.

Any smooth chart on $B(p_j,r_j)$ also gives a smooth chart on
$B(p_s,r_s)$ for all $s > j$. The union of all such charts gives a
smooth atlas on $B_{\infty}$. It is fairly easy to see the metrics
$g_j$ on $B(p_j,r_j)$, converge to a smooth metric $g_{\infty}$ on
$B_{\infty}$ uniformly together with all derivatives on compact
sets. For since the $F_{s-1} \circ \cdots \circ F_j$ are very good
approximate isometries, the $g_j$ are very close to each other,
and hence form a Cauchy sequence (together with their derivatives,
in the sense that the covariant derivatives of $g_j$ with respect
to $g_s$ are very small when $j$ and $s$ are both large). One
checks in the usual way that such a Cauchy sequence converges.

The origins $p_j$ are identified with each other, and hence with
an origin $p_{\infty}$ in $B_{\infty}$. Now it is the inverses of
the maps identifying $B(p_j,r_j)$ with open subsets of
$B_{\infty}$ that provide the diffeomorphisms of (relatively
compact) open sets in $B_{\infty}$ into the geodesic balls
$B(p_j,s_j) \subset M_j$ such that the pull-backs of the metrics
$g_j$ converge to $g_{\infty}$. This completes the proof of Step
3.

\medskip
Now it remains to prove both Claim 1 and Claim 2 in Step 1.

The two claims are proved in the proposition nodes above.

\end{proof}
%$$\eqno\#$$ \vskip 0.3cm

We now want to show how to use this convergence result on
solutions to the Ricci flow. Let us first state the definition for
the convergence of evolving manifolds.

% CZ-BASELINE source-line:8046 source-kind:definition source-id:4.1.3
\begin{definition}[pointed convergence of evolving manifolds]
  \label{def:cz-pointed-convergence-of-evolving-manifolds}
  \dcref{Definition 4.1.3}
  \group{Cao--Zhu Cheeger Type Compactness}
  \source{cao-zhu}
Let $(M_k,g_k(t),p_k)$ be a sequence of evolving marked complete
Riemannian manifolds, with the evolving metrics $g_k(t)$ over a
fixed time interval $t\in(A,\Omega]$, $A < 0 \leq \Omega$, and
with the marked points $p_k\in M_k$. We say a sequence of evolving
marked $(B_0(p_k,s_k),g_k(t),p_k)$ over $t \in (A,\Omega]$, where
$B_0(p_k,s_k)$ are geodesic balls of $(M_k,g_k(0))$ centered at
$p_k$ with the radii $s_k \rightarrow s_{\infty}(\leq +\infty)$,
{\bf converges} in the $C_{\rm loc}^{\infty}$ topology {\bf to an
evolving marked} (maybe noncomplete) {\bf
manifold}\index{converges to an evolving marked manifold}
$(B_{\infty},g_{\infty}(t),p_{\infty})$ over $t\in(A,\Omega]$,
where, at the time $t=0$, $B_{\infty}$ is a geodesic open ball
centered at $p_{\infty}\in B_{\infty}$ with the radius
$s_{\infty}$, if we can find a sequence of exhausting open sets
$U_k$ in $B_{\infty}$ containing $p_{\infty}$ and a sequence of
diffeomorphisms $f_k$ of the sets $U_k$ in $B_{\infty}$ to open
sets $V_k$ in $B(p_k,s_k) \subset M_k$ mapping $p_{\infty}$ to
$p_k$ such that the pull-back metrics
$\tilde{g}_k(t)=(f_k)^*g_k(t)$ converge in $C^\infty$ topology to
$g_{\infty}(t)$ on every compact subset of
$B_{\infty}\times(A,\Omega]$.
\end{definition}

Now we fix a time interval $A<t \leq \Omega$ with $-\infty < A<0$
and $0 \leq \Omega <+\infty$. Consider a sequence of marked
evolving complete manifolds $(M_k,g_k(t),p_k),\ t\in(A,\Omega]$,
with each $g_k(t)$, $k=1, 2, \ldots,$ being a solution of the
Ricci flow
$$
\frac{\partial}{\partial t}g_k(t)=-2\Ric_k(t)
$$
on $B_0(p_k,s_k) \times (A,\Omega]$, where $Ric_k$ is the Ricci
curvature tensor of $g_k$, and  $B_0(p_k,s_k)$ is the geodesic
ball of $(M_k,g_k(0))$ centered at $p_k$ with the radii $s_k
\rightarrow s_{\infty}(\leq +\infty)$.

Assume that for each $r < s_{\infty}$ there are positive constants
$C(r)$ and $k(r)$ such that the curvatures of $g_k(t)$ satisfy the
bound
$$
|Rm(g_k)| \leq C(r)
$$
on $B_0(p_k,r) \times (A,\Omega]$ for all $k \geq k(r).$ We also
assume that $(M_k,g_k(t),p_k)$, $k=1, 2, \ldots,$ have a uniform
injectivity radius bound at the origins $p_k$ at $t=0$. By Shi's
derivatives estimate (Theorem 1.4.1), the above assumption of
uniform bound of the curvatures on the geodesic balls $B_0(p_k,r)$
($r < s_{\infty}$) implies the uniform bounds on all the
derivatives of the curvatures at $t=0$ on the geodesic balls
$B_0(p_k,r)$ ($r < s_{\infty}$). Then by Theorem 4.1.2 we can find
a subsequence of marked evolving manifolds, still denoted by
$(M_k,g_k(t),p_k)$ with $t\in(A,\Omega]$, so that the geodesic
balls $(B_0(p_k,s_k),g_k(0),p_k)$ converge in the $C^\infty_{loc}$
topology to a geodesic ball
$(B_{\infty}(p_{\infty},s_{\infty}),g_{\infty}(0),p_{\infty})$.
>From now on, we consider this subsequence of marked evolving
manifolds.  By Definition 4.1.1, we have a sequence of (relatively
compact) exhausting covering $\{U_k\}$ of
$B_{\infty}(p_{\infty},s_{\infty})$ containing $p_{\infty}$ and a
sequence of diffeomorphisms $f_k$ of the sets $U_k$ in
$B_{\infty}(p_{\infty},s_{\infty})$ to open sets $V_k$ in
$B_{0}(p_{k},s_{k})$ mapping $p_{\infty}$ to $p_k$ such that the
pull-back metrics at $t=0$
$$
\tilde{g}_k(0)=(f_k)^*g_k(0) \stackrel{C_{\rm
loc}^\infty}{\longrightarrow}g_{\infty}(0),\ \quad \text{ as }\;
k\rightarrow+\infty,\; \text{ on }\;
B_{\infty}(p_{\infty},s_{\infty}).
$$
However, the pull-back metrics $\tilde{g}_k(t)=(f_k)^*g_k(t)$ are
also defined at all times $A<t \leq \Omega$ (although
$g_{\infty}(t)$ is not yet). We also have uniform bounds on the
curvature of the pull-back metrics $\tilde{g}_k(t)$ and all their
derivatives, by Shi's derivative estimates (Theorem 1.4.1), on
every compact subset of $B_{\infty}(p_{\infty},s_{\infty}) \times
(A,\Omega]$. What we claim next is that we can find uniform bounds
on all the covariant derivatives of the $\tilde{g}_k$ taken with
respect to the fixed metric $g_{\infty}(0)$.

%\vskip 0.2cm \noindent {\bf Lemma 4.1.4} \emph{\
% CZ-BASELINE source-line:8127 source-kind:lemma source-id:4.1.4
\begin{lemma}[time-uniform smooth convergence from one time slice]
  \label{lem:cz-time-uniform-smooth-convergence-from-one-time-slice}
  \dcref{Lemma 4.1.4}
  \group{Cao--Zhu Cheeger Type Compactness}
  \source{cao-zhu}
Let $(M,g)$ be a Riemannian manifold, $K$ a compact subset of $M$,
and $\tilde{g}_k(t)$ a collection of solutions to Ricci flow
defined on neighborhoods of $K\times[\alpha,\beta]$ with
$[\alpha,\beta]$ containing $0$.  Suppose that for each $l\geq0$,
\begin{itemize}
\item[(a)] $C_0^{-1}g\leq\tilde{g}_k(0)\leq C_0g,\ \ on\ K,\ for\
all\ k,$ \item[(b)] $|\nabla^l\tilde{g}_k(0)|\leq C_l,\ \ on\ K,\
for\ all\ k,$ \item[(c)]
$|\tilde{\nabla}^l_kRm(\tilde{g}_k)|_k\leq C'_l,\ on\
K\times[\alpha,\beta],\ for\ all\ k,$
\end{itemize}
for some positive constants $C_l,\ C'_l,\ l=0,1,\ldots,$
independent of $k$, where $Rm(\tilde{g}_k)$ are the curvature
tensors of the metrics $\tilde{g}_k(t)$, $\tilde{\nabla}_k$ denote
covariant derivative with respect to $\tilde{g}_k(t)$, $|\cdot|_k$
are the length of a tensor with respect to $\tilde{g}_k(t)$, and
$|\cdot|$ is the length with respect to $g$. Then the metrics
$\tilde{g}_k(t)$ satisfy
$$
\tilde{C_0}^{-1}g\leq\tilde{g}_k(t) \leq \tilde{C_0}g,\ \ on\
K\times[\alpha,\beta]
$$
and
$$
|\nabla^l\tilde{g}_k|\leq\tilde{C_l},\ on\ K\times[\alpha,\beta],\
\ l=1,2,\ldots,
$$
for all $k$, where $\tilde{C_l},\ l=0,1,\ldots,$ are positive
constants independent of $k$.
\end{lemma}

%\vskip 0.1cm \noindent {\bf Proof.}
\begin{proof}
First by using the equation
$$
\frac{\partial}{\partial t}\tilde{g}_k=-2\tilde{\Ric}_k
$$
and the assumption (c) we immediately get \be
\tilde{C_0}^{-1}g\leq\tilde{g}_k(t)
\leq \tilde{C_0}g, \; \mbox{ on }\; K\times[\alpha,\beta] %\eqno(4.1.10)
\ee for some positive constant $\tilde{C_0}$ independent of $k$.

Next we want to bound $\nabla\tilde{g}_k$. The difference of the
connection $\tilde{\Gamma}_k$ of $\tilde{g}_k$ and the connection
$\Gamma$ of $g$ is a tensor. Taking $\Gamma$ to be fixed in time,
we get
\begin{align*}
\frac{\partial}{\partial t}(\tilde{\Gamma}_k-\Gamma)& =
\frac{\partial}{\partial
t}\(\frac{1}{2}(\tilde{g}_k)^{\gamma\delta}\left[\frac{\partial}{\partial
x^\alpha}(\tilde{g}_k)_{\delta\beta}+\frac{\partial}{\partial
x^\beta}(\tilde{g}_k)_{\delta\alpha}-\frac{\partial}{\partial
x^\delta}(\tilde{g}_k)_{\alpha\beta}\right]\)\\
& = \frac{1}{2}(\tilde{g}_k)^{\gamma\delta}
\Big[(\tilde{\nabla}_k)_\alpha(-2(\tilde{\Ric}_k)_{\beta\delta})
+(\tilde{\nabla}_k)_\beta(-2(\tilde{\Ric}_k)_{\alpha\delta}) \\
&\qquad-(\tilde{\nabla}_k)_\delta(-2(\tilde{\Ric}_k)_{\alpha\beta})\Big]
\end{align*}
and then by the assumption (c) and (4.1.10),
$$
\left|\frac{\partial}{\partial
t}(\tilde{\Gamma}_k-\Gamma)\right|\leq C,\quad \text{for all }\;
k.
$$
Note also that at a normal coordinate of the metric $g$ at a fixed
point and at the time $t=0$,
\begin{align}
(\tilde{\Gamma}_k)^\gamma_{\alpha\beta}-\Gamma^\gamma_{\alpha\beta}
&=
\frac{1}{2}(\tilde{g}_k)^{\gamma\delta}\(\frac{\partial}{\partial
x^\alpha}(\tilde{g}_k)_{\delta\beta}+\frac{\partial}{\partial
x^\beta}(\tilde{g}_k)_{\delta\alpha}-\frac{\partial}{\partial
x^\delta}(\tilde{g}_k)_{\alpha\beta}\)\\
& =\frac{1}{2}(\tilde{g}_k)^{\gamma\delta}
(\nabla_\alpha(\tilde{g}_k)_{\delta\beta}
+\nabla_\beta(\tilde{g}_k)_{\delta\alpha}
-\nabla_\delta(\tilde{g}_k)_{\alpha\beta}),\nn
\end{align}  %\eqno(4.1.11)
thus by the assumption (b) and (4.1.10),
$$
|\tilde{\Gamma}_k(0)-\Gamma|\leq C, \; \text{ for all }\; k.
$$
Integrating over time we deduce that \be
|\tilde{\Gamma}_k-\Gamma|\leq C, \; \text{ on }\;
K\times[\alpha,\beta],\; \text{ for all }\; k.  %\eqno(4.1.12)
\ee By using the assumption (c) and (4.1.10) again, we have
\begin{align*}
\left|\frac{\partial}{\partial t}(\nabla\tilde{g}_k)\right|
&  = |-2\nabla\tilde{\Ric}_k|\\
&  = |-2\tilde{\nabla}_k\tilde{\Ric}_k+(\tilde{\Gamma}_k
-\Gamma)*\tilde{\Ric}_k|\\
&  \leq C,\; \text{ for all }\; k.
\end{align*}
Hence by combining with the assumption (b) we get bounds \be
|\nabla\tilde{g}_k|\leq\tilde{C}_1,\; \text{ on }\;
K\times[\alpha,\beta],%\eqno(4.1.13)
\ee for some positive constant $\tilde{C}_1$ independent of $k$.

Further we want to bound $\nabla^2\tilde{g}_k$. Again regarding
$\nabla$ as fixed in time, we see
$$
\frac{\partial}{\partial
t}(\nabla^2\tilde{g}_k)=-2\nabla^2(\tilde{R}ic_k).
$$
Write
\begin{align*}
\nabla^2\tilde{\Ric}_k& =
(\nabla-\tilde{\nabla}_k)(\nabla\tilde{\Ric}_k)+
\tilde{\nabla}_k(\nabla-\tilde{\nabla}_k)\tilde{\Ric}_k
+\tilde{\nabla}^2_k\tilde{\Ric}_k\\
&  = (\Gamma-\tilde{\Gamma}_k)*\nabla\tilde{\Ric}_k+
\tilde{\nabla}_k((\Gamma-\tilde{\Gamma}_k)*\tilde{\Ric}_k)
+\tilde{\nabla}^2_k\tilde{\Ric}_k\\
&
=(\Gamma-\tilde{\Gamma}_k)*[(\nabla-\tilde{\nabla}_k)\tilde{\Ric}_k
+\tilde{\nabla}_k\tilde{\Ric}_k] \\
&\quad+
\tilde{\nabla}_k(\tilde{g}_k^{-1}*\nabla\tilde{g}_k*\tilde{\Ric}_k)
+\tilde{\nabla}^2_k\tilde{\Ric}_k\\
&
=(\Gamma-\tilde{\Gamma}_k)*[(\Gamma-\tilde{\Gamma}_k)*\tilde{\Ric}_k
+\tilde{\nabla}_k\tilde{\Ric}_k] \\
&\quad+
\tilde{\nabla}_k(\tilde{g}_k^{-1}*\nabla\tilde{g}_k*\tilde{\Ric}_k)
+\tilde{\nabla}^2_k\tilde{\Ric}_k
\end{align*}
where we have used (4.1.11). Then by the assumption (c), (4.1.10),
(4.1.12) and (4.1.13) we have
\begin{align*}
|\frac{\partial}{\partial t}\nabla^2\tilde{g}_k|
&  \leq  C+C\cdot|\tilde{\nabla}_k\nabla\tilde{g}_k|\\
&  = C+C\cdot|\nabla^2\tilde{g}_k
+(\tilde{\Gamma}_k-\Gamma)*\nabla\tilde{g}_k|\\
&  \leq C+C|\nabla^2\tilde{g}_k|.
\end{align*}
Hence by combining with the assumption (b) we get
$$
|\nabla^2\tilde{g}_k|\leq\tilde{C}_2,\; \mbox{ on }\;
K\times[\alpha,\beta],
$$
for some positive constant $\tilde{C}_2$ independent of $k$.

The bounds on the higher derivatives can be derived by the same
argument. Therefore we have completed the proof of the lemma.
%$$\eqno \#$$ \vskip 0.3cm
\end{proof}

We now apply the lemma to the pull-back metrics
$\tilde{g}_k(t)=(f_k)^*g_k(t)$ on
$B_{\infty}(p_{\infty},s_{\infty}) \times (A,\Omega]$. Since the
metrics $\tilde{g}_k(0)$ have uniform bounds on their curvature
and all derivatives of their curvature on every compact set of
$B_{\infty}(p_{\infty},s_{\infty})$ and converge to the metric
$g_{\infty}(0)$ in $C^\infty_{loc}$ topology, the assumptions (a)
and (b) are certainly held for every compact subset $K \subset
B_{\infty}(p_{\infty},s_{\infty})$ with $g=g_{\infty}(0)$. For
every compact subinterval $[\alpha,\beta] \subset (A,\Omega]$, we
have already seen from Shi's derivative estimates (Theorem 1.4.1)
that the assumption (c) is also held on $K \times [\alpha,\beta]$.
Then all of the $\nabla^l\tilde{g}_k$ are uniformly bounded with
respect to the fixed metric $g = g_{\infty}(0)$ on every compact
set of $B_{\infty}(p_{\infty},s_{\infty})\times(A,\Omega]$. By
using the classical Arzela-Ascoli theorem, we can find a
subsequence which converges uniformly together with all its
derivatives on every compact subset of
$B_{\infty}(p_{\infty},s_{\infty})\times(A,\Omega]$. The limit
metric will agree with that obtained previously at $t=0$, where we
know its convergence already. The limit $g_{\infty}(t),\
t\in(A,\Omega]$, is now clearly itself a solution of the Ricci
flow. Thus we obtain the following Cheeger type compactness
theorem to the Ricci flow, which is essentially obtained by
Hamilton in \cite{Ha95} and is called \textbf{Hamilton's
compactness theorem}\index{Hamilton's compactness theorem}.

%\vskip 0.2cm \noindent {\bf Theorem 4.1.5}

% CZ-BASELINE source-line:8304 source-kind:theorem source-id:4.1.5
\begin{theorem}[Hamilton compactness for Ricci flows]
  \label{thm:cz-hamilton-compactness-for-ricci-flows}
  \dcref{Theorem 4.1.5}
  \group{Cao--Zhu Cheeger Type Compactness}
  \source{cao-zhu}
Let $(M_k,g_k(t),p_k),$ $t\in(A,\Omega]$ with $A<0 \leq \Omega$,
be a sequence of evolving marked complete Riemannian manifolds.
Consider a sequence of geodesic balls $B_0(p_k,s_k) \subset M_k$
of radii $s_k (0<s_k \leq +{\infty})$, with $s_k \rightarrow
s_{\infty}$ $(\leq +{\infty})$, around the base points $p_k$ in
the metrics $g_k(0)$. Suppose each $g_k(t)$ is a solution to the
Ricci flow on $B_0(p_k,s_k) \times (A,\Omega]$. Suppose also
\begin{itemize}
\item[(i)] for every radius $r<s_{\infty}$ there exist positive
constants $C(r)$ and $k(r)$ independent of $k$ such that the
curvature tensors $Rm(g_k)$ of the evolving metrics $g_k(t)$
satisfy the bound
$$
|Rm(g_k)|\leq C(r),
$$
on $B_0(p_k,r) \times (A,\Omega]$ for all $k \geq k(r)$, and
\item[(ii)] there exists a constant $\delta > 0$ such that the
injectivity radii of $M_k$ at $p_k$ in the metric $g_k(0)$ satisfy
the bound
$$
{\rm inj}\,(M_k,p_k,g_k(0))\geq\delta>0,
$$
for all $k = 1, 2, \ldots$.
\end{itemize}
Then there exists a subsequence of evolving marked
$(B_0(p_k,s_k),g_k(t),p_k)$ over $t \in (A,\Omega]$ which converge
in $C^\infty_{loc}$ topology to a solution
$(B_{\infty},g_{\infty}(t),p_{\infty})$ over $t \in (A,\Omega]$ to
the Ricci flow, where, at the time $t=0$, $B_{\infty}$ is a
geodesic open ball centered at $p_{\infty}\in B_{\infty}$ with the
radius $s_{\infty}$. Moreover the limiting solution is complete if
$s_{\infty}=+\infty$.
\end{theorem}
%$$\eqno \#$$ \vskip 1cm


\section{Injectivity Radius Estimates}

We will use rescaling arguments to understand the formation of
singularities and long-time behaviors of the Ricci flow. In view
of the compactness property obtained in the previous section, on
one hand one needs to control the bounds on the curvature, and on
the other hand one needs to control the lower bounds of the
injectivity radius. In applications we usually rescale the
solution so that the (rescaled) curvatures become uniformly
bounded on compact subsets and leave the injectivity radii of the
(rescaled) solutions to be estimated in terms of curvatures. In
this section we will review a number of such injectivity radius
estimates in Riemannian geometry. The combination of these
injectivity estimates with Perelman's no local collapsing theorem
I$'$ yields the well-known little loop lemma to the Ricci flow
which was conjectured by Hamilton in \cite{Ha95F}.

Let $M$ be a Riemannian manifold. Recall that the {\bf injectivity
radius}\index{injectivity radius} at a point $p\in M$ is defined
by
$$
{\rm inj}\,(M,p)=\sup\{r>0\ |\  \exp_p: B(O,r)(\subset T_pM)\to M
\; \text{ is injective}\},
$$
and the {\bf injectivity radius of} ${\mathbf M}$ is
$$
{\rm inj}\,(M)=\inf\{ {\rm inj}\,(M,p)\ |\  p\in M\}.
$$
We begin with a basic lemma due to Klingenberg (cf. Corollary 5.7
of Cheeger \& Ebin \cite{CE}).

\medskip
{\bf Klingenberg's Lemma.}\index{Klingenberg's lemma} \emph{ \ Let
$M$ be a complete Riemannian manifold and let $p\in M$. Let
$l_M(p)$ denote the minimal length of a nontrivial geodesic loop
starting and ending at $p$ $($maybe not smooth at $p)$. Then the
injectivity radius of $M$ at $p$ satisfies the inequality
$$
{\rm inj}\,(M,p)\geq
\min\left\{\frac{\pi}{\sqrt{K_{\max}}},\frac{1}{2}l_M(p)\right\}
$$
where $K_{\max}$ denotes the supermum of the sectional curvature
on $M$ and we understand ${\pi}/{\sqrt{K_{\max}}}$ to be positive
infinity if $K_{\max}\leq0$.}
%$$\eqno \#$$ \vskip 0.3cm

\medskip
Based on this lemma and a second variation argument, Klingenberg
proved that the injectivity radius of an even-dimensional,
compact, simply connected Riemannian manifold of positive
sectional curvature is bounded from below by
$\pi/\sqrt{K_{\max}}$. For odd-dimensional, compact, simply
connected Riemannian manifold of positive sectional curvature, the
same injectivity radius estimates was also proved by Klingenberg
under an additional assumption that the sectional curvature is
strictly $\frac{1}{4}$-pinched (cf. Theorem 5.9 and 5.10 of
\cite{CE}). We also remark that in dimension 7, there exists a
sequence of simply connected, homogeneous Einstein spaces whose
sectional curvatures are positive and uniformly bounded from above
but their injectivity radii converge to zero. (See \cite{AW}.)

The next result, due to Gromoll and Meyer \cite{GM}, shows that
for complete noncompact Riemannian manifold with positive
sectional curvature, the above injectivity radius estimate
actually holds without any restriction on dimensions. Since the
result and proof were not explicitly given in \cite{GM}, we
include a proof here.

%\vskip 2mm \noindent {\bf Theorem 4.2.1 }(\textbf{

% CZ-BASELINE source-line:8411 source-kind:theorem source-id:4.2.1
\begin{theorem}[Gromoll–Meyer injectivity-radius estimate]
  \label{thm:cz-gromollmeyer-injectivity-radius-estimate}
  \dcref{Theorem 4.2.1}
  \group{Cao--Zhu Injectivity Radius Estimates}
  \source{cao-zhu}
\index{Gromoll-Meyer injectivity radius estimate} Let $M$
be a complete, noncompact Riemannian manifold with positive
sectional curvature. Then the injectivity radius of $M$ satisfies
the following estimate
$$
{\rm inj}\,(M)\geq\frac{\pi}{\sqrt{K_{\max}}}.
$$
\end{theorem}

%\vskip 0.1cm \noindent{\bf Proof.}
\begin{proof}
Let $O$ be an arbitrary fixed point in $M$. We need to show that
the injectivity radius at $O$ is not less than
$\pi/\sqrt{K_{\max}}$. We argue by contradiction.  Suppose not,
then by Klingenberg's lemma there exists a closed geodesic loop
$\gamma$ on $M$ starting and ending at $O$ (may be not smooth at
$O$).

Since $M$ has positive sectional curvature, we know from the work
of Gromoll-Meyer \cite{GM} (also cf. Proposition 8.5 of \cite{CE})
that there exists a compact totally convex subset $C$ of $M$
containing the geodesic loop $\gamma$. Among all geodesic loops
starting and ending at the same point and lying entirely in the
compact totally convex set $C$ there will be a shortest one. Call
it $\gamma_0$, and suppose $\gamma_0$ starts and ends at a point
we call $p_0$.

First we claim that $\gamma_0$ must be also smooth at the point
$p_0$. Indeed by the curvature bound and implicit function
theorem, there will be a geodesic loop $\tilde{\gamma}$ close to
$\gamma_0$ starting and ending at any point $\tilde{p}$ close to
$p_0$. Let $\tilde{p}$ be along $\gamma_0$. Then by total
convexity of the set $C$, $\tilde{\gamma}$ also lies entirely in
$C$. If $\gamma_0$ makes an angle different from $\pi$ at $p_0$,
the first variation formula will imply that $\tilde{\gamma}$ is
shorter than $\gamma_0$. This contradicts with the choice of the
geodesic loop $\gamma_0$ being the shortest.

Now let $L:[0,+\infty)\rightarrow M$ be a ray emanating from
$p_0$. Choose $r>0$ large enough and set $q=L(r)$. Consider the
distance between $q$ and the geodesic loop $\gamma_0$. It is clear
that the distance can be realized by a geodesic $\beta$ connecting
the point $q$ to a point $p$ on $\gamma_0$.

Let $X$ be the unit tangent vector of the geodesic loop $\gamma_0$
at $p$. Clearly $X$ is orthogonal to the tangent vector of $\beta$
at $p$. We then translate the vector $X$ along the geodesic
$\beta$ to get a parallel vector field $X(t),\ 0\leq t\leq r$. By
using this vector field we can form a variation fixing one
endpoint $q$ and the other on $\gamma_0$ such that the variational
vector field is $(1-\frac{t}{r})X(t)$. The second variation of the
arclength of this family of curves is given by
\begin{align*}
&  I\(\(1-\frac{t}{r}\)X(t),\(1-\frac{t}{r}\)X(t)\)\\
&  = \int_0^r\bigg[\left|\nabla_{\frac{\partial}{\partial
t}}\(\(1-\frac{t}{r}\)X(t)\)\right|^2 \\
&\qquad-R\(\frac{\partial}{\partial t},\(1-\frac{t}{r}\)X(t),
\frac{\partial}{\partial t},
\(1-\frac{t}{r}\)X(t)\)\bigg]dt\\
&
=\frac{1}{r}-\int_0^r\(1-\frac{t}{r}\)^2R\(\frac{\partial}{\partial
t},X(t), \frac{\partial}{\partial t},X(t)\)dt\\
&  < 0
\end{align*}
when $r$ is sufficiently large, since the sectional curvature of
$M$ is strictly positive everywhere. This contradicts with the
fact that $\beta$ is the shortest geodesic connecting the point
$q$ to the shortest geodesic loop $\gamma_0$. Thus we have proved
the injectivity radius estimate.
%$$\eqno \#$$ \vskip 0.3cm
\end{proof}

In contrast to the above injectivity radius estimates, the
following well-known injectivity radius estimate of Cheeger (cf.
Theorem 5.8 of \cite{CE}) does not impose the restriction on the
sign of the sectional curvature.

\medskip
{\bf Cheeger's Lemma.} \emph{ \ Let $M$ be an $n$-dimensional
compact Riemannian manifold with the sectional curvature
$|K_M|\leq\lambda$, the diameter $d(M)\leq D$, and the volume
$\Vol(M)\geq v > 0$. Then, we have
$$
{\rm inj}\,(M)\ge C_n(\lambda,D,v)
$$
for some positive constant $C_n(\lambda,D,v)$ depending only on
$\lambda,D,v$ and the dimension $n$.}\index{Cheeger's lemma}
%$$\eqno \#$$ \vskip 0.3cm

\medskip

For general manifolds, there are localized versions of the above
Cheeger's Lemma. In 1981, Cheng-Li-Yau \cite{CLY} first obtained
the important estimate of local injectivity radius under the
normalization of the injectivity radius at any fixed base point.
Their result is what Hamilton needed to prove his compactness
result in \cite{Ha95}. Here, we also need the following version of
the local injectivity radius estimate, which appeared in the 1982
paper of Cheeger-Gromov-Taylor \cite{CGT}, in which the
normalization is in terms of local volume of a ball. (According to
Yau, the argument between local volume and local injectivity
radius was, however, initiated by him during a conversation with
Gromov in 1975 in an explanation of his paper \cite{Y76} on
proving complete manifolds with positive Ricci curvature have
infinite volume.)


%\vskip 0.3cm \noindent {\bf Theorem 4.2.2 }(
% CZ-BASELINE source-line:8520 source-kind:theorem source-id:4.2.2
\begin{theorem}[local injectivity-radius estimate from volume and curvature]
  \label{thm:cz-local-injectivity-radius-estimate-from-volume-and-curvature}
  \dcref{Theorem 4.2.2}
  \group{Cao--Zhu Injectivity Radius Estimates}
  \source{cao-zhu}
Let $B(x_0, 4r_0)$, $0<r_0<\infty$, be a geodesic ball in an
$n$-dimensional complete Riemannian manifold $(M,g)$ such that the
sectional curvature $K$ of the metric $g$ on $B(x_0,4r_0)$
satisfies the bounds
$$
\lambda\leq K\leq\Lambda
$$
for some constants $\lambda$ and $\Lambda$. Then for any positive
constant $r\le r_0$ $($we will also require $r\le
\pi/(4\sqrt{\Lambda})$ if $\Lambda>0)$ the injectivity radius of
$M$ at $x_0$ can be bounded from below by
$$
\mbox{\rm inj}(M, x_0)\geq r\cdot\frac{\Vol(B(x_0,
r))}{\Vol(B(x_0,r))+V^n_{\lambda}(2r)},
$$
where $V_{\lambda}^n(2r)$ denotes the volume of a geodesic ball of
radius $2r$ in the $n$-dimensional simply connected space form
$M_{\lambda}$ with constant sectional curvature $\lambda$.
\end{theorem}

%\vskip 0.2cm\noindent {\bf Proof.} \
\begin{proof} The following proof is essentially from Cheeger-Gromov-Taylor \cite{CGT}
(cf. also \cite{AM}).

It is well known (cf. Lemma 5.6 of \cite{CE}) that
$$
\mbox{inj}(M, x_0)=\min \left\{\mbox{conjugate radius of}\ x_0, \
\frac {1}{2}l_M(x_0)\right\}
$$
where $l_M(x_0)$ denotes the length of the shortest (nontrivial)
closed geodesic starting and ending at $x_0$. Since by assumption
$r\le \pi/(4\sqrt{\Lambda})$ if $\Lambda>0$, the conjugate radius
of $x_0$ is at least $4r$. Thus it suffices to show \be
l_M(x_0)\geq
2r\cdot\frac{\Vol(B(x_0,r))}{\Vol(B(x_0,r))+V^n_{\lambda}(2r)}.
%\eqno(4.2.1)
\ee The idea of Cheeger-Gromov-Taylor \cite{CGT} for proving this
inequality, as indicated in \cite{AM}, is to compare the geometry of
the ball $B(x_0,4r)$ $\subseteq B(x_0, 4r_0)\subset M$ with the
geometry of its lifting $\tilde B_{4r}\subset T_{x_0}(M)$, via the
exponential map $\mbox{exp}_{x_0}$, equipped with the pull-back
metric $\tilde g=\mbox{exp}^{*}_{x_0} g$. Thus $\mbox{exp}_{x_0}:
\tilde B_{4r} \to B(x_0,4r)$ is a length-preserving local
diffeomorphism.

Let $\tilde x_0, \tilde x_1, \ldots, \tilde x_N$ be the preimages
of $x_0$ in $\tilde B_r\subset \tilde B_{4r}$ with $\tilde x_0=0$.
Clearly they one-to-one correspond to the geodesic loops
$\gamma_0, \gamma_1, \ldots, \gamma_N$ at $x_0$ of length less
than $r$, where $\gamma_0$ is the trivial loop. Now for each point
$\tilde x_i$ there exists exactly one isometric immersion
$\varphi_i : \tilde B_r \to \tilde B_{4r}$ mapping $0$ to $\tilde
x_i$ and such that $\mbox{exp}_{x_0} \varphi_i=\mbox{exp}_{x_0}$.

Without loss of generality, we may assume $\gamma_1$ is the
shortest nontrivial geodesic loop at $x_0$. By analyzing short
homotopies, one finds that $\varphi_i(\tilde x)\neq
\varphi_j(\tilde x)$ for all $\tilde x\in \tilde B_r$ and $0\le
i<j\le N$. This fact has two consequences:

\medskip
(a) \ $N\ge 2m$, where $m=[r/l_M(x_0)]$. To see this, we first
observe that the points $\varphi_1^k(0), -m\le k\le m$, are
preimages of $x_0$ in $\tilde B_r$ because $\varphi_1$ is an
isometric immersion satisfying $\mbox{exp}_{x_0}
\varphi_1=\mbox{exp}_{x_0}$. Moreover we claim they are distinct.
For otherwise $\varphi_1$ would act as a permutation on the set
$\{\varphi_1^k(0)\ |\ -m\le k\le m\}$. Since the induced metric
$\tilde g$ at each point in $\tilde B_r$ has the injectivity
radius at least $2r$, it follows from the Whitehead theorem (see
for example \cite{CE}) that $\tilde B_r$ is geodesically convex.
Then there would exist the unique center of mass $\tilde y \in
\tilde{B_r}$. But then $\tilde y=\varphi_0(\tilde
y)=\varphi_1(\tilde y)$, a contradiction.

\medskip\noindent
\;\;\;(b) Each point in $B(x_0,r)$ has at least $N\!+\!1$
preimages in $\Omega\!=\!\cup_{i=0}^N B(\tilde x_i, r)$ $\subset
\tilde B_{2r}$. Hence by the Bishop volume comparison,
$$
(N+1)\mbox{Vol}\,(B(x_0,r)) \leq \mbox{Vol}_{\tilde g}(\Omega)
\leq \mbox{Vol}_{\tilde g}(\tilde B_{2r})\le V^n_{\lambda}(2r).
$$

Now the inequality (4.2.1) follows by combining the fact $N\ge
2[r/l_M(x_0)]$ with the above volume estimate.
%$$\eqno \#$$ \vskip 0.3cm
\end{proof}

For our later applications, we now consider in a complete
Riemannian manifold $M$ a geodesic ball $B(p_0, s_0)$ $(0<s_0\le
\infty)$ with the property that there exists a positive increasing
function $\Lambda : [0, s_0)\to [0, \infty)$ such that for any
$0<s<s_0$ the sectional curvature $K$ on the ball $B(p_0, s)$ of
radius $s$ around $p_0$ satisfies the bound
$$
|K|\le \Lambda(s).
$$
Using Theorem 4.2.2, we can control the injectivity radius at any
point $p\in B(p_0, s_0)$ in terms a positive constant that depends
only on the dimension $n$, the injectivity radius at the base
point $p_0$, the function $\Lambda$ and the distance $d(p_0, p)$
from $p$ to $p_0$. We now proceed to derive such an estimate. The
geometric insight of the following argument belongs to Yau
\cite{Y76} where he obtained a lower bound estimate for volume by
comparing various geodesic balls.  Indeed, it is a finite version
of Yau's Busemann function argument which gives the information on
comparing geodesic balls with centers far apart.

For any point $p\in B(p_0, s_0)$ with $d(p_0, p)=s$, set
$r_0=(s_0-s)/4$ (we define $r_0$=1 if $s_0=\infty$). Define the
set $S$ to be the union of minimal geodesic segments that connect
$p$ to each point in $B(p_0, r_0)$. Now any point $q\in S$ has
distance at most
$$
r_0+r_0+s=s+2r_0
$$
from $p_0$ and hence $S\subseteq B(p_0, s+2r_0)$. For any $0<r\le
\min\{\pi/4\sqrt{\Lambda(s+2r_0)},$ $r_0\}$, we denote by $\alpha
(p, r)$ the sector $S\cap B(p,r)$ of radius $r$ and by $\alpha (p,
s+r_0)=S \cap B(p,s+r_0)$. Let $\alpha_{-\Lambda(s+2r_0)}(r_0)$
(resp. $\alpha_{-\Lambda(s+2r_0)} (s+r_0)$) be a corresponding
sector of the same ``angles" with radius $r_0$ (resp. $s+r_0$) in
the $n$-dimensional simply connected space form with constant
sectional curvature $-\Lambda(s+2r_0)$. Since $B(p_0, r_0)\subset
S\subset\alpha (p,s+r_0)$ and $\alpha (p, r)\subset B(p,r)$, the
Bishop-Gromov volume comparison theorem implies that
\begin{align*}
\frac{\Vol(B(p_0,r_0))}{\Vol(B(p,r))}
&\le \frac{\Vol(\alpha (p,s+r_0))}{\Vol(\alpha (p,r))} \\
&\le\frac{\Vol(\alpha_{-\Lambda(s+2r_0)}(s+r_0))}{\Vol(\alpha_
{-\Lambda(s+2r_0)}(r))}=\frac{V^n_{-\Lambda(s+2r_0)}(s+r_0)}{V^n_
{-\Lambda(s+2r_0)}(r)}.
\end{align*}
Combining this inequality with the local injectivity radius
estimate in Theorem 4.2.2, we get
\begin{align*}
&{\rm inj}\,(M,p) \\
&\geq r \frac{V^{n}_{-\Lambda(s+2r_0)}(r)\cdot
\Vol(B(p_0,r_0))}{V^{n}_{-\Lambda(s+2r_0)}(r)\Vol(B(p_0,r_0))+
V_{-\Lambda(s+2r_0)}^n(2r)V^{n}_{-\Lambda(s+2r_0)}(s+2r_0)}.
\end{align*}
Thus, we have proved the following (cf. \cite{CLY}, \cite{CGT} or
\cite{AM})

%\vskip 0.2cm\noindent {\bf Corollary 4.2.3} \ \emph{
% CZ-BASELINE source-line:8667 source-kind:corollary source-id:4.2.3
\begin{corollary}[injectivity radius from local curvature scale]
  \label{cor:cz-injectivity-radius-from-local-curvature-scale}
  \dcref{Corollary 4.2.3}
  \group{Cao--Zhu Injectivity Radius Estimates}
  \source{cao-zhu}
Suppose $B(p_0, s_0)$ $(0<s_0\le \infty)$ is a geodesic ball in an
$n$-dimensional complete Riemannian manifold $M$ having the
property that for any $0<s<s_0$ the sectional curvature $K$ on
$B(p_0, s)$ satisfies the bound
$$
|K|\le \Lambda(s)
$$
for\; some\; positive\; increasing\; function\; $\Lambda$\;
defined\; on\; $[0,\,s_0)$.\; Then\; for\; any\; point\; $p\;\in\;
B(p_0,\, s_0)$\; with\; $d(p_0,\, p)\;=\;s$\; and\; any\;
positive\; number $r\le \min\{\pi/4\sqrt{\Lambda(s+2r_0)},r_0\}$
with $r_0=(s_0-s)/4$, the injectivity radius of $M$ at $p$ is
bounded below by
\begin{align*}
& {\rm inj}\,(M,p) \\
&\geq r \frac{V^{n}_{-\Lambda(s+2r_0)}(r)\cdot
\Vol(B(p_0,r_0))}{V^{n}_{-\Lambda(s+2r_0)}(r)\Vol(B(p_0,r_0))+
V_{-\Lambda(s+2r_0)}^n(2r)V^{n}_{-\Lambda(s+2r_0)}(s+2r_0)}.
\end{align*}
In particular, we have \be
\inj(M,p)\geq \rho_{n, \delta, \Lambda}(s) %\eqno(4.2.2)
\ee where $\delta>0$ is a lower bound of the injectivity radius
$\inj(M, p_0)$ at the origin $p_0$ and $\rho_{n,\delta, \Lambda}:
[0, s_0) \to \mathbb{R^+}$ is a positive decreasing function that
depends only on the dimension $n$, the lower bound $\delta$ of the
injectivity radius $\inj(M,p_0)$, and the function $\Lambda$.
%$$\eqno \#$$ \vskip 0.3cm
\end{corollary}

We remark that in the above discussion if $s_0=\infty$ then we can
apply the standard Bishop relative volume comparison theorem to
geodesic balls directly. Indeed, for any $p\in M$ and any positive
constants $r$ and $r_0$, we have $B(p_0, r_0)\subseteq B(p,
\hat{r})$ with
$\hat{r}\stackrel{\Delta}{=}\max\{r,r_0+d(p_0,p)\}$. Suppose in
addition the curvature $K$ on $M$ is uniformly bounded by
$\lambda\le K\le\Lambda$ for some constants $\lambda$ and
$\Lambda$, then the Bishop volume comparison theorem implies that
$$
\frac{\Vol(B(p_0,r_0))}{\Vol(B(p,r))}\le
\frac{\Vol(B(p,\hat{r}))}{\Vol(B(p,r))}\le
\frac{V_{\lambda}(\hat{r})}{V_{\lambda}(r)}.$$ Hence \be
\inj(M,p)\geq r\frac{V_{\lambda}^n(r)\cdot
\Vol(B(p_0,r_0))}{V_{\lambda}^n(r)\Vol(B(p_0,r_0))
+V_{\lambda}^n(2r)V_{\lambda}^n(\hat{r})}. %\eqno(4.2.3)
\ee So we see that the injectivity radius $inj(M,p)$ at $p$ falls
off at worst exponentially as the distance $d(p_0,p)$ goes to
infinity. In other words, \be \inj(M,p)\geq
\frac{c}{\sqrt{B}}(\delta \sqrt{B})^n
e^{-C\sqrt{B}d(p,p_0)} %\eqno(4.2.5) ----> (4.2.4)---CAO
\ee where $B$ is an upper bound on the absolute value of the
sectional curvature, $\delta$ is a lower bound on the injectivity
radius at $p_0$ with $\delta<c/\sqrt{B}$, and $c>0$ and $C<+\infty
$ are positive constants depending only on the dimension $n$.

\smallskip
Finally, by combining Theorem 4.2.2 with Perelman's no local
collapsing Theorem I${'}$ (Theorem 3.3.3) we immediately obtain the
following important (due to Perelman \cite{P1}) \textbf{Little Loop
Lemma} \index{Little Loop Lemma} conjectured by Hamilton
\cite{Ha95F}.

%\vskip 0.2cm \noindent {\bf Theorem 4.2.4}
% CZ-BASELINE source-line:8731 source-kind:theorem source-id:4.2.4
\begin{theorem}[little loop lemma]
  \label{thm:cz-little-loop-lemma}
  \dcref{Theorem 4.2.4}
  \group{Cao--Zhu Injectivity Radius Estimates}
  \source{cao-zhu}
Let $g_{ij}(t)$, $0\le t<T<+\infty$, be a solution of the Ricci
flow on a compact manifold $M$. Then there exists a constant
$\rho>0$ having the following property: if at a point $x_0\in M$
and a time $t_0\in [0,T)$,
$$ |Rm|(\cdot,t_0)\leq r^{-2} \quad
\text{on}\ B_{t_0}(x_0,r)
$$
for some $r\leq \sqrt{T}$, then the injectivity radius of $M$ with
respect to the metric $g_{ij}(t_0)$ at $x_0$ is bounded from below
by
$$
\inj(M,x_0,g_{ij}(t_0))\geq \rho r.
$$
%$$\eqno \#$$
\end{theorem}


\section{Limiting Singularity Models}

In \cite{Ha95F}, Hamilton classified singularities of the Ricci
flow into three types and showed each type has a corresponding
singularity model. The main purpose of this section is to discuss
these results of Hamilton. Most of the presentation follows
Hamilton \cite{Ha95F}.

Consider a solution $g_{ij}(x,t)$ of the Ricci flow on $M\times
[0,T)$, $T\le +\infty$, where either $M$ is compact or at each
time $t$ the metric $g_{ij}(\cdot,t)$ is complete and has bounded
curvature. We say that $g_{ij}(x,t)$ is a {\bf maximal solution}
\index{maximal solution} if either $T=+\infty$ or $T<+\infty$ and
$|Rm|$ is unbounded as $t\to T$.

Denote by
$$
K_{\max}(t)=\sup_{x\in M}|Rm(x,t)|_{g_{ij}(t)}.
$$

%\vskip 0.2cm \noindent{\bf Definition 4.3.1}\ \
% CZ-BASELINE source-line:8770 source-kind:definition source-id:4.3.1
\begin{definition}[almost maximum curvature points]
  \label{def:cz-almost-maximum-curvature-points}
  \dcref{Definition 4.3.1}
  \group{Cao--Zhu Limiting Singularity Models}
  \source{cao-zhu}
We say that $\{(x_k,t_k) \in M\times [0,T) \}$, $k=1, 2, \ldots$,
is a sequence of {\bf (almost) maximum points}\index{(almost)
maximum points} if there exist positive constants $c_1$ and
$\alpha \in (0,1]$ such that
$$
|Rm(x_k,t_k)|\geq c_1 K_{\max}(t), \quad t\in
[t_k-\frac{\alpha}{K_{\max}(t_k)},t_k]
$$
for all $k$.
\end{definition}

%\vskip 0.3cm \noindent {\bf Definition 4.3.2} \
% CZ-BASELINE source-line:8783 source-kind:definition source-id:4.3.2
\begin{definition}[injectivity-radius condition at almost maximum points]
  \label{def:cz-injectivity-radius-condition-at-almost-maximum-points}
  \dcref{Definition 4.3.2}
  \group{Cao--Zhu Limiting Singularity Models}
  \source{cao-zhu}
We say that the solution satisfies {\bf injectivity radius
condition}\index{injectivity radius!condition} if for any sequence
of (almost) maximum points $\{(x_k,t_k)\}$, there exists a
constant $c_2>0$ independent of $k$ such that
$$
\inj(M,x_k,g_{ij}(t_k))\geq \frac{c_2}{\sqrt{K_{\max}(t_k)}} \quad
\text{for all }\;   k.
$$
\end{definition}

Clearly, by the Little Loop Lemma, a maximal solution on a compact
manifold with the maximal time $T<+\infty$ always satisfies the
injectivity radius condition. Also by the Gromoll-Meyer
injectivity radius estimate, a solution on a complete noncompact
manifold with positive sectional curvature also satisfies the
injectivity radius condition.

According to Hamilton \cite{Ha95F}, we can classify maximal
solutions into three types; every maximal solution is clearly of
one and only one of the following three types:
$$
\mbox{\textbf{Type I:}} \ \ \ T<+\infty\ \  \mbox{and}\ \
\sup_{t\in [0,T)}(T-t)K_{\max}(t)<+\infty;\
$$
\index{type I}
$$
\ \ \ \mbox{\textbf{Type II:}} \ \ \ \text{(a)}\  T<+\infty\ \
\mbox{but}\ \ \sup_{t\in [0,T)}(T-t)K_{\max}(t)=+\infty;
$$
\index{type II!(a)}
$$
 \hskip 2.3cm \text{(b)}\  T=+\infty\ \  \mbox{but}\ \
\sup_{t\in [0,T)}tK_{\max}(t)=+\infty;\ \ \ \ \ \
$$
\index{type II!(b)}
$$
\mbox{\textbf{Type III:}} \ \ \ \text{(a)}\  T=+\infty,\ \ \ \
\sup_{t\in [0,T)}tK_{\max}(t)<+\infty,\ \ \mbox{and} \ \
$$
$$ \limsup_{t\rightarrow
+\infty}tK_{\max}(t)>0;
$$
\index{type III!(a)}
$$
\hskip 2.3cm \text{(b)}\  T=+\infty,\ \ \ \  \sup_{t\in
[0,T)}tK_{\max}(t)<+\infty,\ \ \mbox{and} \ \
$$
$$ \limsup_{t\rightarrow
+\infty}tK_{\max}(t)=0;
$$
\index{type III!(b)}

Note that Type III (b) is not compatible with the injectivity
radius condition unless it is a trivial flat solution. (Indeed
under the Ricci flow the length of a curve $\gamma$ connecting two
points $x_0$, $x_1\in M$ evolves by
\begin{align*}
\frac{d}{dt}L_t(\gamma)
&  = \int_{\gamma}-\Ric(\dot{\gamma},\dot{\gamma})ds\\
&  \leq C(n)K_{\max}(t)\cdot L_t(\gamma)\\
&  \leq \frac{\epsilon}{t}L_t(\gamma),\; \text{ as $t$ large
enough,}
\end{align*}
for arbitrarily fixed $\epsilon>0$. Thus when we are considering
the Ricci flow on a compact manifold, the diameter of the evolving
manifold grows at most as $t^{\epsilon}$. But the curvature of the
evolving manifold decays faster than $t^{-1}$. This says, as
choosing $\epsilon>0$ small enough,
$$
{\rm diam}_t(M)^2\cdot |Rm(\cdot,t)|\rightarrow0,\; \text{ as }\;
t\rightarrow+\infty.
$$
Then it is well-known from Cheeger-Gromov \cite{Gro78} that the
manifold is a nilmanifold and the injectivity radius condition can
not be satisfied as $t$ large enough. When we are considering the
Ricci flow on a complete noncompact manifold with nonnegative
curvature operator or on a complete noncompact K\"ahler manifold
with nonnegative holomorphic bisectional curvature,
Li-Yau-Hamilton inequalities imply that $tR(x,t)$ is increasing in
time $t$. Then Type III(b) occurs only when the solution is a
trivial flat metric.)

For each type of maximal solution, Hamilton \cite{Ha95F} defined a
corresponding type of limiting singularity model.

%\vskip 0.2cm\noindent {\bf Definition 4.3.3} \ \
% CZ-BASELINE source-line:8870 source-kind:definition source-id:4.3.3
\begin{definition}[singularity types for maximal Ricci flows]
  \label{def:cz-singularity-types-for-maximal-ricci-flows}
  \dcref{Definition 4.3.3}
  \group{Cao--Zhu Limiting Singularity Models}
  \source{cao-zhu}
A solution $g_{ij}(x,t)$ to the Ricci flow on the manifold $M$,
where either $M$ is compact or at each time $t$ the metric
$g_{ij}(\cdot,t)$ is complete and has bounded curvature, is called
a {\bf singularity model}\index{singularity model} if it is not
flat and of one of the following three types:

\vskip 0.1cm\noindent \textbf{Type I}\index{type I}: The solution
exists for $t\in(-\infty,\Omega)$ for some constant $\Omega$ with
$0<\Omega<+\infty$ and
$$
|Rm|\leq\Omega/(\Omega-t)
$$
everywhere with equality somewhere at $t=0$;

\vskip 0.1cm\noindent \textbf{Type II}\index{type II}: The
solution exists for $t\in(-\infty,+\infty)$  and
$$
|Rm|\leq 1 \ \ \ \ \ \ \ \ \
$$
everywhere with equality somewhere at $t=0$;

\vskip 0.1cm\noindent \textbf{Type III}\index{type III}: The
solution exists for $t\in(-A,+\infty)$ for some constant $A$ with
$0<A<+\infty$ and
$$
\ \ |Rm|\leq A/(A+t)
$$
everywhere with equality somewhere at $t=0$.
\end{definition}

%\vskip 0.3cm \noindent {\bf Theorem 4.3.4} \ \ \emph{
% CZ-BASELINE source-line:8902 source-kind:theorem source-id:4.3.4
\begin{theorem}[existence of dilation-limit singularity models]
  \label{thm:cz-existence-of-dilation-limit-singularity-models}
  \dcref{Theorem 4.3.4}
  \group{Cao--Zhu Limiting Singularity Models}
  \source{cao-zhu}
For any maximal solution to the Ricci flow which satisfies the
injectivity radius condition and is of Type {\rm I, II(a), (b),}
or {\rm III(a),} there exists a sequence of dilations of the
solution along (almost) maximum points which converges in the
$C_{\rm loc}^{\infty}$ topology to a singularity model of the
corresponding type.
\end{theorem}

%\vskip 0.1cm \noindent {\bf Proof.}
\begin{proof} \

\smallskip
{\bf Type I:}\index{type I} \ We consider a maximal solution
$g_{ij}(x,t)$ on $M\times [0,T)$ with $T<+\infty$ and
$$
\Omega\stackrel{\Delta}{=}\limsup_{t\rightarrow T}
(T-t)K_{\max}(t)<+\infty.
$$

First we note that $\Omega>0$. Indeed by the evolution equation of
curvature,
$$
\frac{d}{dt}K_{\max}(t)\leq{\rm Const}\, \cdot K_{\max}^2(t).
$$
This implies that
$$
K_{\max}(t)\cdot (T-t)\geq {\rm Const}\, >0,
$$
because
$$
\limsup_{t\rightarrow T}K_{\max}(t)=+\infty.
$$
Thus $\Omega$ must be positive.

Choose a sequence of points $x_k$ and times $t_k$ such that
$t_k\rightarrow T$ and
$$
\lim_{k\rightarrow \infty}(T-t_k)|Rm(x_k,t_k)|=\Omega.
$$
Denote by
$$
\epsilon_k=\frac{1}{\sqrt{|Rm(x_k,t_k)|}}.
$$
We translate in time so that $t_k$ becomes 0, dilate in space by
the factor $\epsilon_k$ and dilate in time by $\epsilon_k^2$ to
get
$$
\tilde{g}_{ij}^{(k)}(\cdot,\tilde{t})
=\epsilon_k^{-2}g_{ij}(\cdot,t_k+\epsilon_k^2\tilde{t}),\ \ \ \
\tilde{t}\in [-t_k/\epsilon_k^2, (T-t_k)/\epsilon_k^2).
$$
Then
\begin{align*}
\frac{\partial}{\partial
\tilde{t}}\tilde{g}_{ij}^{(k)}(\cdot,\tilde{t})& =
\epsilon_k^{-2}\frac{\partial
}{\partial t}g_{ij}(\cdot,t)\cdot \epsilon_k^2 \\
&  = -2R_{ij}(\cdot,t_k+\epsilon_k^2\tilde{t}) \\
&  = -2\tilde{R}_{ij}^{(k)}(\cdot,\tilde{t}),
\end{align*}
where $\tilde{R}_{ij}^{(k)}$ is the Ricci curvature of the metric
$\tilde{g}_{ij}^{(k)}$. So $\tilde{g}_{ij}^{(k)}(\cdot,\tilde{t})$
is still a solution to the Ricci flow which exists on the time
interval $[-t_k/\epsilon_k^2,(T-t_k)/\epsilon_k^2)$, where
$$
t_k/\epsilon_k^2=t_k|Rm(x_k,t_k)|\to +\infty
$$
and
$$
(T-t_k)/\epsilon_k^2=(T-t_k)|Rm(x_k,t_k)|\to \Omega.
$$

For any $\epsilon>0$ we can find a time $\tau <T $ such that for
$t\in [\tau,T)$,
$$
|Rm|\leq (\Omega+\epsilon)/(T-t)
$$
by the assumption. Then for $\tilde{t}\in
[(\tau-t_k)/\epsilon_k^2,(T-t_k)/\epsilon_k^2)$, the curvature of
$\tilde{g}_{ij}^{(k)}(\cdot,\tilde{t})$ is bounded by
\begin{align*}
|\tilde{R}m^{(k)}|&  = \epsilon_k^2|Rm| \\
&  \leq(\Omega+\epsilon)/((T-t)|Rm(x_k,t_k)|) \\
& =(\Omega+\epsilon)/((T-t_k)|Rm(x_k,t_k)|+(t_k-t)|Rm(x_k,t_k)|)\\
&  \rightarrow (\Omega+\epsilon)/(\Omega-\tilde{t}), \qquad
\text{as } \; k\rightarrow +\infty.
\end{align*}
This implies that $\{(x_k,t_k)\}$ is a sequence of (almost)
maximum points. And then by the injectivity radius condition and
Hamilton's compactness theorem 4.1.5, there exists a subsequence
of the metrics $\tilde{g}_{ij}^{(k)}(\tilde {t})$ which converges
in the $C_{loc}^{\infty}$ topology to a limit metric
$\tilde{g}_{ij}^{(\infty)}(\tilde {t})$ on a limiting manifold
$\tilde{M}$ with $\tilde{t}\in(-\infty,\Omega)$ such that
$\tilde{g}_{ij}^{(\infty)}(\tilde {t})$ is a complete solution of
the Ricci flow and its curvature satisfies the bound
$$
|\tilde{R}m^{(\infty)}|\leq \Omega/(\Omega-\tilde{t})
$$
everywhere on $\tilde{M}\times (-\infty,\Omega)$ with the equality
somewhere at $\tilde{t}=0$.

\medskip
{\bf Type II(a):}\index{type II!(a)} \ We consider a maximal
solution $g_{ij}(x,t)$ on $M\times [0,T)$ with
$$
T<+\infty \quad\text{and} \quad \limsup_{t\rightarrow
T}(T-t)K_{\max}(t)=+\infty.
$$

Let $T_k<T<+\infty$ with $T_k\rightarrow T$, and $\gamma_k\nearrow
1$, as $k\rightarrow +\infty$. Pick points $x_k$ and times $t_k$
such that, as $k\rightarrow +\infty$,
$$
(T_k-t_k)|Rm(x_k,t_k)|\geq \gamma_k \sup_{x\in M, t\leq
T_k}(T_k-t)|Rm(x,t)|\rightarrow +\infty.
$$
Again denote by
$$
\epsilon_k=\frac{1}{\sqrt{|Rm(x_k,t_k)|}}
$$
and dilate the solution as before to get
$$
\tilde{g}_{ij}^{(k)}(\cdot,\tilde{t})
=\epsilon_k^{-2}g_{ij}(\cdot,t_k+\epsilon_k^2\tilde{t}),\ \ \ \
\tilde{t}\in [-t_k/\epsilon_k^2,(T_k-t_k)/\epsilon_k^2),
$$
which is still a solution to the Ricci flow and satisfies the
curvature bound
\begin{align*}
|\tilde{R}m^{(k)}|&  = \epsilon_k^2|Rm| \\
&  \leq\frac{1}{\gamma_k}\cdot \frac{(T_k-t_k)}{(T_k-t)}\\
& =\frac{1}{\gamma_k}\frac{(T_k-t_k)|Rm(x_k,t_k)|}{[(T_k-t_k)|
Rm(x_k,t_k)|-\tilde{t}]}\quad \mbox{ for } \; \tilde{t}\in\
\bigg[-\frac{t_k}{\epsilon_k^2},\frac{(T_k-t_k)}{\epsilon_k^2}\bigg),
\end{align*}
since $t=t_k+\epsilon_k^2\tilde{t}$ and
$\epsilon_k=1/\sqrt{|Rm(x_k,t_k)|}$. Hence $\{(x_k,t_k)\}$ is a
sequence of (almost) maximum points. And then as before, by
applying Hamilton's compactness theorem 4.1.5, there exists a
subsequence of the metrics $\tilde{g}_{ij}^{(k)}(\tilde{t})$ which
converges in the $C_{loc}^{\infty}$ topology to a limit
$\tilde{g}_{ij}^{(\infty)}(\tilde {t})$ on a limiting manifold
$\tilde{M}$ and $\tilde{t}\in(-\infty,+\infty)$ such that
$\tilde{g}_{ij}^{(\infty)}(\tilde {t})$ is a complete solution of
the Ricci flow and its curvature satisfies
$$
|\tilde{R}m^{(\infty)}|\leq 1
$$
everywhere on $\tilde{M}\times (-\infty,+\infty)$ and the equality
holds somewhere at $\tilde{t}=0$.

\medskip
{\bf Type II(b):}\index{type II!(b)}\ \ We consider a maximal
solution $g_{ij}(x,t)$ on $M\times [0,T)$ with
$$
T=+\infty\quad \text{and} \quad \limsup_{t\rightarrow
T}tK_{\max}(t)=+\infty.
$$

Again let $T_k\rightarrow T=+\infty$, and $\gamma_k\nearrow 1$, as
$k\rightarrow +\infty$. Pick $x_k$ and $t_k$ such that
$$
t_k(T_k-t_k)|Rm(x_k,t_k)|\geq \gamma_k \sup_{x\in M, t\leq
T_k}t(T_k-t)|Rm(x,t)|.
$$
Define
$$
\tilde{g}_{ij}^{(k)}(\cdot,\tilde{t})
=\epsilon_k^{-2}g_{ij}(\cdot,t_k+\epsilon_k^2\tilde{t}),\ \ \ \
\tilde{t}\in [-t_k/\epsilon_k^2,(T_k-t_k)/\epsilon_k^2),
$$
where $\epsilon_k=1/\sqrt{|Rm(x_k,t_k)|}$.\vskip 0.1cm \noindent
Since
\begin{align*}
t_k(T_k-t_k)|Rm(x_k,t_k)|&  \geq \gamma_k \sup\limits_{x\in M,
t\leq T_k}t(T_k-t)|Rm(x,t)|\\
&  \geq\gamma_k \sup\limits_{x\in M, t\leq T_k/2}t(T_k-t)|Rm(x,t)|\\
&  \geq\frac{T_k}{2}\gamma_k \sup\limits_{x\in M, t\leq
T_k/2}t|Rm(x,t)|,
\end{align*}
we have
$$
\frac{t_k}{\epsilon_k^2}=t_k|Rm(x_k,t_k)|\geq
\frac{\gamma_k}{2}\(\frac{T_k}{T_k-t_k}\)\sup_{x\in M, t\leq
T_k/2}t|Rm(x,t)| \to +\infty,
$$
and
$$
\frac{(T_k-t_k)}{\epsilon_k^2}=(T_k-t_k)|Rm(x_k,t_k)|\geq
\frac{\gamma_k}{2}\(\frac{T_k}{t_k}\)\sup_{x\in M, t\leq
T_k/2}t|Rm(x,t)| \to +\infty,
$$
as $k \rightarrow +\infty$. As before, we also have
$$
\frac{\partial}{\partial
\tilde{t}}\tilde{g}_{ij}^{(k)}(\cdot,\tilde{t})
=-2\tilde{R}_{ij}^{(k)}(\cdot,\tilde{t})
$$
and {\small
\begin{align*}
& |\tilde{R}m^{(k)}| \\
&  =\epsilon_k^2|Rm| \\
&  \leq \frac{1}{\gamma_k} \cdot \frac{t_k(T_k-t_k)}{t(T_k-t)}\\
&  =\frac{1}{\gamma_k}\cdot \frac{t_k(T_k-t_k)|Rm(x_k,t_k)|}
{(t_k+\epsilon_k^2\tilde{t})
[(T_k-t_k)-\epsilon_k^2\tilde{t}]\cdot |Rm(x_k,t_k)|}\\
&  =\frac{1}{\gamma_k}\cdot \frac{t_k(T_k-t_k)|Rm(x_k,t_k)|}
{(t_k+\epsilon_k^2\tilde{t})[(T_k-t_k)|Rm(x_k,t_k)|-\tilde{t}]}\\
&  =\frac{t_k(T_k-t_k)|Rm(x_k,t_k)|}{\gamma_k
(\!1\!+\!\tilde{t}/(t_k|Rm(x_k,t_k)|))[t_k(T_k\!-\!t_k)|Rm(x_k,t_k)|]
(1\!-\!\tilde{t}/((T_k\!-\!t_k)|Rm(x_k,t_k)|))}\\
&  \rightarrow 1,\quad \text{as } \; k\rightarrow+\infty.
\end{align*}
} Hence $\{(x_k,t_k)\}$ is again a sequence of (almost) maximum
points. As before, there exists a subsequence of the metrics
$\tilde{g}_{ij}^{(k)}(\tilde {t})$ which converges in the
$C_{loc}^{\infty}$ topology to a limit
$\tilde{g}_{ij}^{(\infty)}(\tilde {t})$ on a limiting manifold
$\tilde{M}$ and $\tilde{t}\in(-\infty,+\infty)$ such that
$\tilde{g}_{ij}^{(\infty)}(\tilde {t})$ is a complete solution of
the Ricci flow and its curvature satisfies
$$
|\tilde{R}m^{(\infty)}|\leq 1
$$
everywhere on $\tilde{M}\times (-\infty,+\infty)$ with the
equality somewhere at $\tilde{t}=0$.

\medskip
{\bf Type III(a):}\index{type III!(a)} \ We consider a maximal
solution $g_{ij}(x,t)$ on $M\times [0,T)$ with $T=+\infty$ and
$$
\limsup\limits_{t\rightarrow T}tK_{\max}(t)=A\in (0,+\infty).
$$
Choose a sequence of $x_k$ and $t_k$ such that
$t_k\rightarrow+\infty$ and
$$
\lim_{k\rightarrow \infty}t_k|Rm(x_k,t_k)|=A.
$$
Set $\epsilon_k=1/\sqrt{|Rm(x_k,t_k)|}$ and dilate the solution as
before to get
$$
\tilde{g}_{ij}^{(k)}(\cdot,\tilde{t})
=\epsilon_k^{-2}g_{ij}(\cdot,t_k+\epsilon_k^2\tilde{t}),\ \ \ \
\tilde{t}\in [-t_k/\epsilon_k^2,+\infty)
$$
which is still a solution to the Ricci flow. Also, for arbitrarily
fixed $\epsilon>0$, there exists a sufficiently large positive
constant $\tau$ such that for $t\in [\tau,+\infty)$,
\begin{align*}
|\tilde{R}m^{(k)}|&  = \epsilon_k^2|Rm| \\
&  \leq \epsilon_k^2\(\frac{A+\epsilon}{t}\)\\
&  = \epsilon_k^2 \(\frac{A+\epsilon}{t_k+\epsilon_k^2\tilde{t}}\)\\
&  = (A+\epsilon)/(t_k|Rm(x_k,t_k)|+\tilde{t}),\quad \text{for }\;
\tilde{t}\in [(\tau-t_k)/\epsilon_k^2,+\infty).
\end{align*}
Note that
$$
(A+\epsilon)/(t_k|Rm(x_k,t_k)|+\tilde{t})\rightarrow
(A+\epsilon)/(A+\tilde{t}), \quad \text{as } \; k\rightarrow
+\infty
$$
and
$$
(\tau-t_k)/\epsilon_k^2\rightarrow -A,\quad \text{as }\;
k\rightarrow +\infty.
$$
Hence $\{(x_k,t_k)\}$ is a sequence of (almost) maximum points.
And then as before, there exists a subsequence of the metrics
$\tilde{g}_{ij}^{(k)}(\tilde {t})$ which converges in the
$C_{loc}^{\infty}$ topology to a limit
$\tilde{g}_{ij}^{(\infty)}(\tilde {t})$ on a limiting manifold
$\tilde{M}$ and $\tilde{t}\in(-A,+\infty)$ such that
$\tilde{g}_{ij}^{(\infty)}(\tilde {t})$ is a complete solution of
the Ricci flow and its curvature satisfies
$$
|\tilde{R}m^{(\infty)}| \leq A/(A+\tilde{t})
$$
everywhere on $\tilde{M}\times (-A,+\infty)$ with  the equality
somewhere at $\tilde{t}=0$.
%$$\eqno \#$$
\end{proof}

In the case of manifolds with nonnegative curvature operator, or
K\"ahler metrics with nonnegative holomorphic bisectional curvature,
we can bound the Riemannian curvature by the scalar curvature $R$ up
to a constant factor depending only on the dimension. Then we can
slightly modify the statements in the previous theorem as follows

%\vskip 0.2cm \noindent{\bf Corollary 4.3.5} \ \emph{

% CZ-BASELINE source-line:9194 source-kind:corollary source-id:4.3.5
\begin{corollary}[restricted singularity-model types under nonnegative curvature]
  \label{cor:cz-restricted-singularity-model-types-under-nonnegative-curvature}
  \dcref{Corollary 4.3.5}
  \group{Cao--Zhu Limiting Singularity Models}
  \source{cao-zhu}
For any complete maximal solution to the Ricci flow satisfying the
injectivity radius condition and with bounded and nonnegative
curvature operator on a Riemannian manifold, or on a K\"ahler
manifold with bounded and nonnegative holomorphic bisectional
curvature, there exists a sequence of dilations of the solution
along $($almost$\,)$ maximum points which converges to a singular
model.

\vskip 0.1cm \noindent For Type {\rm I} solutions:\  the limit
model exists for $t\in (-\infty,\Omega)$ with $0<\Omega<+\infty$
and has
$$
R\leq \Omega/(\Omega-t)
$$
everywhere with equality somewhere at $t=0$.

\vskip 0.1cm \noindent For Type {\rm II} solutions:\ the limit
model exists for $t\in (-\infty, +\infty)$  and has
$$
R\leq 1
$$
everywhere with equality somewhere at $t=0$.

\vskip 0.1cm \noindent For Type {\rm III} solutions:\ the limit
model exists for $t\in(-A, +\infty)$ with $0<A<+\infty$ and has
$$
R\leq A/(A+t)
$$
everywhere with equality somewhere at $t=0$.
%$$\eqno\#$$\vskip 0.3cm
\end{corollary}

A natural and important question is to understand each of the three
types of singularity models. The following results obtained by
Hamilton \cite{Ha93E} and Chen-Zhu \cite{CZ00} characterize the Type
II and Type III singularity models with nonnegative curvature
operator and positive Ricci curvature respectively.

%\vskip 0.2cm \noindent {\bf Theorem 4.3.6}
% CZ-BASELINE source-line:9234 source-kind:theorem source-id:4.3.6
\begin{theorem}[Type II models with positive Ricci curvature are steady solitons]
  \label{thm:cz-type-ii-models-with-positive-ricci-curvature-are-steady-solitons}
  \dcref{Theorem 4.3.6}
  \group{Cao--Zhu Limiting Singularity Models}
  \source{cao-zhu}
\
\begin{itemize}
\item[(i)] \ {\rm (Hamilton \cite{Ha93E})} \ Any Type {\rm II}
singularity model with nonnegative curvature operator and positive
Ricci curvature to the Ricci flow on a manifold $M$ must be a
$($steady$)$ Ricci soliton. \item[(ii)] \ {\rm (Chen-Zhu
\cite{CZ00})} \ Any Type {\rm III} singularity model with
nonnegative curvature operator and positive Ricci curvature on a
manifold $M$ must be a homothetically expanding Ricci soliton.
\end{itemize}
\end{theorem}

%\vskip 0.1cm \noindent {\bf Proof.}\ \
\begin{proof}
  \uses{thm:cz-matrix-liyauhamilton-estimate}
We first prove (i).  Translate the time coordinate of the eternal Type II
model by an arbitrary constant $A>0$ and apply the Li--Yau--Hamilton
inequality on the interval beginning at time $-A$.  Letting $A\to\infty$
removes every $1/t$ term and gives
\[
\frac{\partial R}{\partial t}+2\nabla_iR\,V^i+2R_{ij}V^iV^j\ge0.
\]
At a spacetime point where $R$ assumes its maximum, this quadratic
vanishes for $V=0$.  The strong maximum argument carried out below for
(ii), now with all $1/t$ terms deleted, propagates a unique smooth null
vector field $V$.  Its first variation satisfies
$\nabla_iR+2R_{ij}V^j=0$.  Equations (4.3.7)--(4.3.14), with the same
deletion of the $1/t$ terms, then give
\[
4R^{ij}g^{kl}(R_{ik}-\nabla_kV_i)(R_{jl}-\nabla_kV_j)=0.
\]
Positive definiteness of the Ricci tensor implies
$\nabla_iV_j=R_{ij}$.  Equivalently, the metric evolves only by the
one-parameter family of diffeomorphisms generated by $V$, so it is a
steady Ricci soliton.

We now prove (ii).

After a shift of the time variable, we may assume the Type III
singularity model is defined on $0<t<+\infty$ and $tR$ assumes its
maximum in space-time.

Recall from the Li-Yau-Hamilton inequality (Theorem 2.5.4) that
for any vectors $V^i$ and $W^i$, \be
M_{ij}W^iW^j+(P_{kij}+P_{kji})V^kW^iW^j+R_{ikjl}W^iW^jV^kV^l
\geq 0, %(4.3.1)
\ee where
$$
M_{ij}=\Delta R_{ij}-\frac{1}{2}\nabla_i\nabla_jR+2R_{ipjq}R^{pq}
-g^{pq}R_{ip}R_{jq}+\frac{1}{2t}R_{ij}
$$
and
$$
P_{ijk}=\nabla_iR_{jk}-\nabla_jR_{ik}.
$$
Take the trace on $W$ to get \be
Q\stackrel{\Delta}{=}\frac{\partial R}{\partial
t}+\frac{R}{t}+2\nabla_iR\cdot V^i+2R_{ij}V^iV^j\geq0
%\eqno(4.3.2)
\ee for any vector $V^i$. Let us choose $V$ to be the vector field
minimizing $Q$, i.e., \be
V^i=-\frac{1}{2}(\Ric^{-1})^{ik}\nabla_kR, %\eqno(4.3.3)
\ee where $(Ric^{-1})^{ik}$ is the inverse of the Ricci tensor
$R_{ij}$. Substitute this vector field $V$ into $Q$ to get a smooth
function $\tilde{Q}$. The calculations of Chow-Hamilton in the proof
of Theorem 6.1 of \cite{CH} give, \be \frac{\partial}{\partial
t}\tilde{Q}
\ge \Delta \tilde{Q}-\frac{2}{t}\tilde{Q}. %\eqno(4.3.4)
\ee

Suppose $tR$ assumes its maximum at $(x_0,t_0)$ with $t_0>0$, then
$$
\frac{\partial R}{\partial t}+\frac{R}{t}=0,\qquad \text{at}\quad
(x_0,t_0).
$$
This implies that the quantity
$$
Q=\frac{\partial R}{\partial t}+\frac{R}{t} +2\nabla_iR\cdot
V^i+2R_{ij}V^iV^j
$$
vanishes in the direction $V=0$ at $(x_0,t_0)$. We claim that for
any earlier time $t<t_0$ and any point $x\in M$, there is a vector
$V\in T_xM$ such that $Q=0$.

We argue by contradiction. Suppose not, then there is $\bar{x}\in
M$ and $0<\bar{t}<t_0$ such that $\tilde{Q}$ is positive at $x =
\bar{x}$ and $t=\bar{t}$. We can find a nonnegative smooth
function $\rho$ on $M$ with support in a neighborhood of $\bar{x}$
so that $\rho(\bar{x})>0$ and
$$
\tilde{Q}\geq \frac{\rho}{\bar{t}^2},
$$
at $t=\bar{t}$. Let $\rho$ evolve by the heat equation
$$
\frac{\partial\rho}{\partial t}=\Delta \rho.
$$
It then follows from the standard strong maximum principle that
$\rho>0$ everywhere for any $t>\bar{t}$. From (4.3.4) we see that
$$
\frac{\partial}{\partial t}\(\tilde{Q}-\frac{\rho}{t^2}\)
\ge\Delta\(\tilde{Q}-\frac{\rho}{t^2}\)-\frac{2}{t}\(\tilde{Q}
-\frac{\rho}{t^2}\)
$$
Then by the maximum principle as in Chapter 2, we get
$$
\tilde{Q}\ge \frac{\rho}{t^2}>0,\qquad \text{for all }\; t\ge
\bar{t}.
$$
This gives a contradiction with the fact $Q=0$ for $V=0$ at
$(x_0,t_0)$. We thus prove the claim.

Consider each time $t<t_0$. The null vector field of $Q$ satisfies
the equation \be
\nabla_iR+2R_{ij}V^j=0,%\eqno(4.3.5)
\ee by the first variation of $Q$ in $V$. Since $R_{ij}$ is
positive, we see that such a null vector field is unique and
varies smoothly in space-time.

Substituting (4.3.5) into the expression of $Q$, we have \be
\frac{\partial R}{\partial t}+\frac{R}{t}
+\nabla_iR\cdot V^i=0 %\eqno(4.3.6)
\ee Denote by
$$
Q_{ij}=M_{ij}+(P_{kij}+P_{kji})V^k+R_{ikjl}V^kV^l.
$$

{}From (4.3.1) we see that $Q_{ij}$ is nonnegative definite with
its trace $Q=0$ for such a null vector $V$. It follows that
$$
Q_{ij}=M_{ij}+(P_{kij}+P_{kji})V^k+R_{ikjl}V^kV^l=0.
$$
Again from the first variation of $Q_{ij}$ in $V$, we see that \be
(P_{kij}+P_{kji})+(R_{ikjl}+R_{jkil})V^l=0, %\eqno(4.3.7)
\ee and hence \be
M_{ij}-R_{ikjl}V^kV^l=0. %\eqno(4.3.8)
\ee

Applying the heat operator to (4.3.5) and (4.3.6) we get
\begin{align}
0&  = \(\frac{\partial}{\partial
t}-\Delta\)(\nabla_iR+2R_{ij}V^j)\\
&  = 2R_{ij}\(\frac{\partial}{\partial
t}-\Delta\)V^j+\(\frac{\partial}{\partial
t}-\Delta\)(\nabla_iR) \nn\\
&\quad+2V^j\(\frac{\partial}{\partial
t}-\Delta\)R_{ij}-4\nabla_kR_{ij}\nabla^kV^j, \nn
\end{align}  %\eqno(4.3.9)
and
\begin{align}
0&  = \(\frac{\partial}{\partial t}-\Delta\)\(\frac{\partial
R}{\partial t}+\frac{R}{t}+\nabla_iR\cdot V^i\)\\
&  = \nabla_iR\(\frac{\partial}{\partial
t}-\Delta\)V^i+V^i\(\frac{\partial}{\partial
t}-\Delta\)(\nabla_iR)-2\nabla_k\nabla_iR\cdot \nabla^kV^i \nn\\
&\quad +\(\frac{\partial}{\partial t}-\Delta\)\(\frac{\partial
R}{\partial t}+\frac{R}{t}\). \nn
\end{align}   %\eqno(4.3.10)

Multiplying (4.3.9) by $V^i$, summing over $i$ and adding
(4.3.10), as well as using the evolution equations on curvature,
we get
\begin{align}
0&  =
2V^i(2\nabla_i(|Rc|^2)-R_{il}\nabla^lR)+2V^iV^j(2R_{piqj}R^{pq}
-2g^{pq}R_{pi}R_{qj})\\
&\quad -4\nabla_kR_{ij}\cdot\nabla^kV^j\cdot V^i-2\nabla_k\nabla_i
R\cdot\nabla^kV^i+4R^{ij}\nabla_i\nabla_jR \nn\\
&\quad +4g^{kl}g^{mn}g^{pq}R_{km}R_{np}R_{ql}
+4R_{ijkl}R^{ik}V^jV^l-\frac{R}{t^2}. \nn
\end{align}  %\eqno(4.3.11)
{}From (4.3.5), we have the following equalities \be
\begin{cases}
-2V^iR_{il}\nabla^lR-4V^iV^jg^{pq}R_{pi}R_{qj}=0,\\
-4\nabla_kR_{ij}\cdot\nabla^kV^j\cdot V^i-2\nabla_k\nabla_i
R\cdot\nabla^kV^i=4R_{ij}\nabla_kV^i\cdot\nabla^kV^j,\\
\nabla_i\nabla_jR=-2\nabla_iR_{jl}\cdot V^l-2R_{jl}\nabla_iV^l.
\end{cases}   %%\eqno(4.3.12)
\ee Substituting (4.3.12) into (4.3.11), we obtain
\begin{align*}
& 8R^{ij}(\nabla_kR_{ij}\cdot
V^k+R_{ikjl}V^kV^l-\nabla_iR_{jl}\cdot
V^l-R_{jl}\nabla_iV^l)\\
&\qquad+4R_{ij}\nabla_kV^i\cdot\nabla^kV^j
+4g^{kl}g^{mn}g^{pq}R_{km}R_{np}R_{ql}-\frac{R}{t^2}=0.
\end{align*}
By using (4.3.7), we know
$$
R^{ij}(\nabla_kR_{ij}\cdot V^k +R_{ikjl}V^kV^l-\nabla_iR_{jl}\cdot
V^l)=0.
$$
Then we have \be
-8R^{ij}R_{jl}\nabla_iV^l+4R_{ij}\nabla_kV^i\cdot\nabla^kV^j
+4g^{kl}g^{mn}g^{pq}R_{km}R_{np}R_{ql}-\frac{R}{t^2}=0.%\eqno(4.3.13)
\ee By taking the trace in the last equality in (4.3.12) and using
(4.3.6) and the evolution equation of the scalar curvature, we can
get \be
R^{ij}(R_{ij}+\frac{g_{ij}}{2t}-\nabla_iV_j)=0. %\eqno(4.3.14)
\ee Finally by combining (4.3.13) and (4.3.14), we deduce
$$
4R^{ij}g^{kl}\(R_{ik}+\frac{g_{ik}}{2t}-\nabla_kV_i\)
\(R_{jk}+\frac{g_{jk}}{2t}-\nabla_kV_j\)=0.
$$
Since $R_{ij}$ is positive definite, we get \be
\nabla_iV_j=R_{ij}+\frac{g_{ij}}{2t},\qquad
\mbox{for all }\;i,j.  %\eqno(4.3.15)
\ee This means that $g_{ij}(t)$ is a homothetically expanding
Ricci soliton.
%$$\eqno \#$$ \vskip 0.3cm
\end{proof}

%\noindent{\bf Remark 4.3.7}\;
% CZ-BASELINE source-line:9424 source-kind:remark source-id:4.3.7
\begin{remark}[noncompact Type II models and steady solitons]
  \label{rem:cz-noncompact-type-ii-models-and-steady-solitons}
  \dcref{Remark 4.3.7}
  \group{Cao--Zhu Limiting Singularity Models}
  \source{cao-zhu}
  \uses{thm:cz-type-ii-models-with-positive-ricci-curvature-are-steady-solitons}
Recall from Section 1.5 that any compact steady Ricci soliton or
expanding Ricci soliton must be Einstein. If the manifold $M$ in
Theorem 4.3.6 is noncompact and simply connected, then the steady
(or expanding) Ricci soliton must be a steady (or expanding)
gradient Ricci soliton. For example, we know that $\nabla_iV_j$ is
symmetric from (4.3.15).  Also, by the simply connectedness of $M$
there exists a function $F$ such that
$$
\nabla_i\nabla_jF=\nabla_iV_j,\qquad \text{on }\;M.
$$
So
$$
R_{ij}=\nabla_i\nabla_jF-\frac{g_{ij}}{2t},\qquad \text{on }\;M
$$
This means that $g_{ij}$ is an expanding gradient Ricci soliton.
\end{remark}

In the K\"ahler case, we have the following results for Type II
and Type III singularity models with nonnegative holomorphic
bisectional curvature obtained by the first author in
\cite{Cao97}.

%\vskip 0.2cm\noindent{\bf Theorem 4.3.8} (
% CZ-BASELINE source-line:9448 source-kind:theorem source-id:4.3.8
\begin{theorem}[Kähler Type II models are steady solitons]
  \label{thm:cz-kahler-type-ii-models-are-steady-solitons}
  \dcref{Theorem 4.3.8}
  \group{Cao--Zhu Limiting Singularity Models}
  \source{cao-zhu}
 \
\begin{itemize}
\item[(i)] Any Type {\rm II} singularity model on a K\"ahler
manifold with nonnegative holomorphic bisectional curvature and
positive Ricci curvature must be a steady K\"ahler-Ricci soliton.
\item[(ii)] Any Type {\rm III} singularity model on a K\"ahler
manifold with nonnegative holomorphic bisectional curvature and
positive Ricci curvature must be an expanding K\"ahler-Ricci
soliton.
\end{itemize}
\end{theorem}
%$$\eqno\#$$ \vskip 0.3cm

To conclude this section, we state a result of Sesum \cite{Se} on
compact Type I singularity models. Recall that Perelman's
functional $\mathcal{W}$, introduced in Section 1.5, is given by
$$
\mathcal{W}(g,f,\tau)=\int_M(4\pi\tau)^{-\frac{n}{2}}[\tau(|\nabla
f|^2+R)+f-n]e^{-f}dV_g
$$
with the function $f$ satisfying the constraint
$$
\int_M(4\pi\tau)^{-\frac{n}{2}}e^{-f}dV_g=1.
$$
And recall from Corollary 1.5.9 that
$$
\mu(g(t))=\inf\left\{\mathcal{W}(g(t),f,T-t)|
\int_M(4\pi(T-t))^{-\frac{n}{2}}e^{-f}dV_{g(t)}=1\right\}
$$
is strictly increasing along the Ricci flow unless we are on a
gradient shrinking soliton. If one can show that $\mu(g(t))$ is
uniformly bounded from above and the minimizing functions
$f=f(\cdot,t)$ have a limit as $t\rightarrow T$, then the
rescaling limit model will be a shrinking gradient soliton. As
shown by Natasa Sesum in \cite{Se}, Type I assumption guarantees
the boundedness of $\mu(g(t))$, while the compactness assumption
of the rescaling limit guarantees the existence of the limit for
the minimizing functions $f(\cdot,t)$. Therefore we have

%\vskip 0.2cm \noindent {\bf Theorem 4.3.9} (
% CZ-BASELINE source-line:9488 source-kind:theorem source-id:4.3.9
\begin{theorem}[compact Type I limits are shrinking solitons]
  \label{thm:cz-compact-type-i-limits-are-shrinking-solitons}
  \dcref{Theorem 4.3.9}
  \group{Cao--Zhu Limiting Singularity Models}
  \source{cao-zhu}
Let $(M,g_{ij}(t))$ be a Type {\rm I} singularity model obtained
as a rescaling limit of a Type {\rm I} maximal solution. Suppose
$M$ is compact. Then $(M,g_{ij}(t))$ must be a gradient shrinking
Ricci soliton.
\end{theorem}
%$$\eqno\#$$\vskip 0.3cm

It seems that the assumption on the compactness of the rescaling
limit is superfluous. We conjecture that any noncompact Type I
limit is also a gradient shrinking soliton.


\section{Ricci Solitons}

We now follow Hamilton \cite{Ha95F} to examine the structure of a
steady Ricci soliton that we get as a Type II limit . The material
is from section 20 of Hamilton \cite{Ha95F} and Hamilton
\cite{Ha88}.

%\vskip 0.2cm \noindent{\bf Lemma 4.4.1}
% CZ-BASELINE source-line:9509 source-kind:lemma source-id:4.4.1
\begin{lemma}[scalar-curvature identities for steady gradient solitons]
  \label{lem:cz-scalar-curvature-identities-for-steady-gradient-solitons}
  \dcref{Lemma 4.4.1}
  \group{Cao--Zhu Ricci Solitons}
  \source{cao-zhu}
Suppose we have a complete gradient steady Ricci soliton $g_{ij}$
with bounded curvature so that
$$
R_{ij}=\nabla_i\nabla_jF
$$
for some function $F$ on $M$. Assume the Ricci curvature is
positive and the scalar curvature $R$ attains its maximum
$R_{\max}$ at a point $x_0\in M$. Then \be
|\nabla F|^2+R=R_{\max} %\eqno(4.4.1)
\ee everywhere on $M$, and furthermore $F$ is convex and attains
its minimum at $x_0$.
\end{lemma}

%\vskip 0.1cm\noindent{\bf Proof.}\;
\begin{proof}
Recall that, from (1.1.15) and noting our $F$ here is $-f$ there,
the steady gradient Ricci soliton has the property
$$
|\nabla F|^2+R=C_0
$$
for some constant $C_0$. Clearly, $C_0\ge R_{\max}$.

If $C_0=R_{\max}$, then $\nabla F=0$ at the point $x_0$. Since
$\nabla_i\nabla_jF=R_{ij}>0,$ we see that $F$ is convex and $F$
attains its minimum at $x_0$.

If $C_0>R_{\max}$, consider a gradient path of $F$ in a local
coordinate neighborhood through $x_0=(x_0^1,\ldots,x_0^n):$
$$
\left\{\arraycolsep=1.5pt\begin{array}{l} x^i=x^i(u),\qquad
u\in(-\varepsilon,\varepsilon),\;i=1,\ldots,n\\[2mm]
x_0^i=x^i(0),
\end{array}\right.
$$
and
$$
\frac{dx^i}{du} =g^{ij}\nabla_jF,\quad
u\in(-\varepsilon,\varepsilon).
$$
Now $|\nabla F|^2=C_0-R\ge C_0-R_{\max}>0$ everywhere, while
$|\nabla F|^2$ is smallest at $x=x_0$ since $R$ is largest there.
But we compute
\begin{align*}
\frac{d}{du}|\nabla
F|^2&  = 2g^{jl}\(\frac{d}{du}\nabla_jF\)\nabla_lF\\
&  = 2g^{ik}g^{jl}\nabla_i\nabla_jF\cdot\nabla_kF\nabla_lF\\
&  = 2g^{ik}g^{jl}R_{ij}\nabla_kF\nabla_lF\\
&  > 0
\end{align*}
since $R_{ij}>0$ and $|\nabla F|^2>0$. Then $|\nabla F|^2$ is not
smallest at $x_0$, and we have a contradiction. %$$\eqno \#$$
\end{proof}

We remark that when we are considering a complete expanding
gradient Ricci soliton on $M$ with positive Ricci curvature and
$$
R_{ij}+\rho g_{ij}=\nabla_i\nabla_jF
$$
for some constant $\rho>0$ and some function $F$, the above
argument gives
$$
|\nabla F|^2+R-2\rho F=C
$$
for some positive constant $C$. Moreover the function $F$ is an
exhausting and convex function. In particular, such an expanding
gradient Ricci soliton is diffeomorphic to the Euclidean space
$\mathbb R^n$.

Let us introduce a geometric invariant as follows. Let $O$ be a
fixed point in a Riemannian manifold $M$, $s$ the distance to the
fixed point $O$, and $R$ the scalar curvature. We define the {\bf
asymptotic scalar curvature ratio}\index{asymptotic scalar
curvature ratio}
$$
A=\limsup_{s\rightarrow +\infty}Rs^2.
$$
Clearly the definition is independent of the choice of the fixed
point $O$ and invariant under dilation. This concept is particular
useful on manifolds with positive sectional curvature. The first
type of gap theorem was obtained by Mok-Siu-Yau \cite{MoSY} in
understanding the hypothesis of the paper of Siu-Yau \cite{SiY}.
Yau (see \cite{GW82}) suggested that this should be a general
phenomenon. This was later conformed by Greene-Wu \cite{GW82, GW},
Eschenberg-Shrader-Strake \cite{ESS} and Drees \cite{Dr} where
they show that any complete noncompact $n$-dimensional (except
$n=4$ or 8) Riemannian manifold of positive sectional curvature
must have $A>0$. Similar results on complete noncompact
K$\rm\ddot{a}$hler manifolds of positive holomorphic bisectional
curvature were obtained by Chen-Zhu \cite{CZ03} and Ni-Tam
\cite{NT}.

%\vskip 0.2cm \noindent{\bf Theorem 4.4.2}\ (
% CZ-BASELINE source-line:9602 source-kind:theorem source-id:4.4.2
\begin{theorem}[asymptotic scalar curvature of positively curved steady solitons]
  \label{thm:cz-asymptotic-scalar-curvature-of-positively-curved-steady-solitons}
  \dcref{Theorem 4.4.2}
  \group{Cao--Zhu Ricci Solitons}
  \source{cao-zhu}
For a complete noncompact steady gradient Ricci soliton with
bounded curvature and positive sectional curvature of dimension
$n\ge3$ where the scalar curvature assume its maximum at a point
$O\in M$, the asymptotic scalar curvature ratio is infinite, i.e.,
$$
A=\limsup_{s\rightarrow +\infty}Rs^2=+\infty
$$
where s is the distance to the point $O$.
\end{theorem}

%\vskip 0.1cm \noindent {\bf Proof.}\
\begin{proof}
The solution to the Ricci flow corresponding to the soliton exists
for $-\infty<t<+\infty$ and is obtained by flowing along the
gradient of a potential function $F$ of the soliton. We argue by
contradiction. Suppose $Rs^2\le C$. We will show that the limit
$$
\bar{g}_{ij}(x)=\lim_{t\rightarrow-\infty}g_{ij}(x,t)
$$
exists for $x\neq O$ on the manifold $M$ and is a complete flat
metric on $M\setminus\{O\}.$ Since the sectional curvature of $M$
is positive everywhere, it follows from Cheeger-Gromoll \cite{CG}
that $M$ is diffeomorphic to $\mathbb{R}^n$. Thus
$M\setminus\{O\}$ is diffeomorphic to
$\mathbb{S}^{n-1}\times\mathbb{R}$. But for $n\ge 3$ there is no
flat metric on $\mathbb{S}^{n-1}\times\mathbb{R}$, and this will
finish the proof.

To see the limit metric exists, we note that $R\rightarrow 0$ as
$s\rightarrow+\infty$, so $|\nabla F|^2\rightarrow R_{\max}$ as
$s\rightarrow +\infty$ by (4.4.1). The function $F$ itself can be
taken to evolve with time, using the definition
$$
\frac{\partial F}{\partial t} =\nabla_iF\cdot\frac{\partial
x^i}{\partial t} =-|\nabla F|^2 =\Delta F-R_{\max}
$$
which pulls $F$ back by the flow along the gradient of $F$. Then
we continue to have $\nabla_i\nabla_jF=R_{ij}$ for all time, and
$|\nabla F|^2\rightarrow R_{\max}$ as $s\rightarrow +\infty$ for
each time.

When we go backward in time, this is equivalent to flowing
outwards along the gradient of $F$, and our speed approaches
$\sqrt{R_{\max}}$. So, starting outside of any neighborhood of $O$
we have
$$
\frac{s}{|t|}=\frac{d_t(\cdot,O)}{|t|}\rightarrow
\sqrt{R_{\max}},\quad \text{as }\;t\rightarrow-\infty
$$
and \be R(\cdot,t)\le\frac{C}{R_{\max}\cdot|t|^2},\quad
\text{as}\;|t|\;\text{large enough}.  %\eqno(4.4.2)
\ee Hence for $|t|$ sufficiently large,
\begin{align*}
0&  \ge -2R_{ij}\\
&  = \frac{\partial}{\partial t}g_{ij}\\
&  \ge -2Rg_{ij}\\
&  \ge -\frac{2C}{R_{\max}\cdot|t|^2}g_{ij}
\end{align*}
which implies that for any tangent vector $V$,
$$
0\le \frac{d}{d|t|}(\log(g_{ij}(t)V^iV^j))
\le\frac{2C}{R_{\max}\cdot|t|^2}.
$$
These two inequalities show that $g_{ij}(t)V^iV^j$ has a limit
$\bar{g}_{ij}V^iV^j$ as $t\rightarrow-\infty$.

Since the metrics are all essentially the same, it always takes an
infinite length to get out to the infinity. This shows the limit
$\bar{g}_{ij}$ is complete at the infinity. One the other hand,
any point $P$ other than $O$ will eventually be arbitrarily far
from $O$, so the limit metric $\bar{g}_{ij}$ is also complete away
from $O$ in $M\setminus\{O\}$. Using Shi's derivative estimates in
Chapter 1, it follows that $g_{ij}(\cdot,t)$ converges in the
$C_{\rm loc}^\infty$ topology to a complete smooth limit metric
$\bar{g}_{ij}$ as $t\rightarrow-\infty$, and the limit metric is
flat by (4.4.2).
%$$\eqno \#$$ \vskip 0.3cm
\end{proof}

The above argument actually shows that \be
\limsup_{s\rightarrow+\infty}Rs^{1+\varepsilon}=+\infty
%\eqno(4.4.3)
\ee for arbitrarily small $\varepsilon>0$ and for any complete
gradient Ricci soliton with bounded and positive sectional
curvature of dimension $n\ge3$ where the scalar curvature assumes
its maximum at a fixed point $O$.

Finally we conclude this section with the important result of
Hamilton on the uniqueness of complete Ricci soliton on
two-dimensional Riemannian manifolds.

%\vskip 0.2cm \noindent {\bf Theorem 4.4.3}\ (
% CZ-BASELINE source-line:9695 source-kind:theorem source-id:4.4.3
\begin{theorem}[classification of two-dimensional steady solitons]
  \label{thm:cz-classification-of-two-dimensional-steady-solitons}
  \dcref{Theorem 4.4.3}
  \group{Cao--Zhu Ricci Solitons}
  \source{cao-zhu}
The only complete steady Ricci soliton on a two-dimensional
manifold with bounded curvature which assumes its maximum $1$ at
an origin is the ``cigar" soliton on the plane $\mathbb{R}^2$ with
the metric
$$
ds^2=\frac{dx^2+dy^2}{1+x^2+y^2}.
$$
\end{theorem}

%\noindent{\bf Proof.}\
\begin{proof}
  \uses{lem:cz-scalar-curvature-identities-for-steady-gradient-solitons}
Recall that the scalar curvature evolves by
$$
\frac{\partial R}{\partial t}=\Delta R+R^2
$$
on a two-dimensional manifold $M$. Denote by
$R_{\min}(t)=\inf\{R(x,t)\ |\ x\in M\}.$ We see from the maximum
principle (see for example Chapter 2) that $R_{\min}(t)$ is
strictly increasing whenever $R_{\min}(t)\neq 0$, for
$-\infty<t<+\infty$. This shows that the curvature of a steady
Ricci soliton on a two-dimensional manifold $M$ must be
nonnegative and $R_{\min}(t) = 0$ for all $t \in
(-\infty,+\infty)$. Further by the strong maximum principle we see
that the curvature is actually positive everywhere. In particular,
the manifold must be noncompact. So the manifold $M$ is
diffiomorphic to $\mathbb{R}^2$ and the Ricci soliton must be a
gradient soliton. Let $F$ be a potential function of the gradient
Ricci soliton. Then, by definition, we have
$$
\nabla_iV_j+\nabla_jV_i=Rg_{ij}
$$
with $V_i=\nabla_iF$. This says that the vector field $V$ must be
conformal. In complex coordinate a conformal vector field is
holomorphic. Hence $V$ is locally given by
$V(z)\frac{\partial}{\partial z}$ for a holomorphic function
$V(z)$. At a zero of $V$ there will be a power series expansion
$$
V(z)=az^p+\cdots,\;(a\neq 0)
$$
and if $p>1$ the vector field will have closed orbits in any
neighborhood of the zero. Now the vector field is gradient and a
gradient flow cannot have a closed orbit. Hence $V(z)$ has only
simple zeros. By Lemma 4.4.1, we know that $F$ is strictly convex
with the only critical point being the minima, chosen to be the
origin of $\mathbb{R}^2$. So the holomorphic vector field $V$ must
be
$$
V(z)\frac{\partial}{\partial z} =cz\frac{\partial}{\partial
z},\qquad \text{for}\;z\in C,
$$
for some complex number $c$.

We now claim that $c$ is real. Let us write the metric as
$$
ds^2=g(x,y)(dx^2+dy^2)
$$
with $z=x+\sqrt{-1}y.$ Then $\nabla F=cz\frac{\partial}{\partial
z}$ means that if $c=a+\sqrt{-1}b$, then
$$
\frac{\partial F}{\partial x}=(ax-by)g,\quad \frac{\partial
F}{\partial y}=(bx+ay)g.
$$
Taking the mixed partial derivatives $\frac{\partial^2 F}{\partial
x\partial y}$ and equating them at the origin $x=y=0$ gives $b=0$,
so $c$ is real.

Let
$$\left\{\arraycolsep=1.5pt\begin{array}{l} x=e^u\cos
v,\qquad -\infty<u<+\infty,\\[2mm]
y=e^u\sin v,\qquad 0\le v\le 2\pi.
\end{array}\right.
$$
Write
\begin{align*}
ds^2& = g(x,y)(dx^2+dy^2)\\
&  = g(e^u\cos v,e^u\sin v)e^{2u}(du^2+dv^2)\\
&  \stackrel{\Delta}{=} g(u,v)(du^2+dv^2).
\end{align*}
Then we get the equations
$$
\frac{\partial F}{\partial u}=ag,\quad \frac{\partial F}{\partial
v}=0
$$
since the gradient of $F$ is just $a\frac{\partial}{\partial u}$
for a real constant $a$. The second equation shows that $F=F(u)$
is a function of $u$ only, then the first equation shows that
$g=g(u)$ is also a function of $u$ only. Then we can write the
metric as
\begin{align}
ds^2&  = g(u)(du^2+dv^2)\\
&  = g(u)e^{-2u}(dx^2+dy^2).\nn
%\eqno(4.4.4)
\end{align}
This implies that $e^{-2u}g(u)$ must be a smooth function of
$x^2+y^2=e^{2u}$. So as $u\rightarrow -\infty$, \be
g(u)=b_1e^{2u}+b_2(e^{2u})^2+\cdots,%\eqno(4.4.5)
\ee with $b_1>0$.

The curvature of the metric is given by
$$
R=-\frac{1}{g}\(\frac{g'}{g}\)'
$$
where $(\cdot)'$ is the derivative with respect to $u$. Note that
the soliton is by translation in $u$ with velocity $c$. Hence
$g=g(u+ct)$ satisfies
$$
\frac{\partial g}{\partial t}=-Rg
$$
which becomes
$$
cg'=\(\frac{g'}{g}\)'.
$$
Thus by (4.4.5),
$$
\frac{g'}{g}=cg+2
$$
and then by integrating
$$
e^{2u}\(\frac{1}{g}\)=-\frac{c}{2}e^{2u}+b_1
$$
i.e.,
$$
g(u)=\frac{e^{2u}}{b_1-\frac{c}{2}e^{2u}}.
$$
In particular, we have $c<0$ since the Ricci soliton is not flat.
Therefore
$$
ds^2=g(u)e^{-2u}(dx^2+dy^2)
=\frac{dx^2+dy^2}{\alpha_1+\alpha_2(x^2+y^2)}
$$
for some constants $\alpha_1,\alpha_2>0$. By the normalization
condition that the curvature attains its maximum $1$ at the
origin, we conclude that
$$
ds^2=\frac{dx^2+dy^2}{1+(x^2+y^2)}.
$$
\end{proof}
%$$\eqno \#$$\vskip 1cm


%\bigskip
